%=======================================================================
%  MONOGRAFÍA / TRABAJO DE GRADO
%  THUNDERS-Master: un proceso ligero para construir, monitorear y
%  validar el entendimiento compartido en ingeniería de requisitos.
%
%  Plantilla lista para Overleaf (compilar con pdfLaTeX).
%  Estructura basada en la monografía de referencia (Agredo-Delgado, 2022).
%
%  CONVENCIONES DE EDICIÓN
%  -----------------------
%  * Las secciones ya redactadas contienen el contenido derivado de:
%       - Revisión sistemática de la literatura (RSL)
%       - Caracterización de elicitación, especificación y validación
%       - Adaptación situacional de THUNDERS (SME / method chunks)
%       - Especificación operativa de THUNDERS-Master
%  * Las secciones pendientes están marcadas con la caja \pendiente{...}
%    (color amarillo). Son las "etapas vacías" para completar a futuro.
%  * Reemplaza los campos de portada (director, codirector, fecha) según
%    corresponda a tu programa.
%=======================================================================
\documentclass[12pt,letterpaper,oneside]{report}

%------------------------- Idioma y codificación -----------------------
\usepackage[utf8]{inputenc}
\usepackage[T1]{fontenc}
\usepackage[spanish,es-tabla]{babel}   % en Overleaf el español funciona sin problema
\usepackage{lmodern}
\usepackage{amssymb}

%------------------------------ Geometría ------------------------------
\usepackage[letterpaper,top=2.5cm,bottom=2.5cm,left=3cm,right=2.5cm]{geometry}
\usepackage{setspace}
\onehalfspacing

%------------------------------ Gráficos -------------------------------
\usepackage{graphicx}
\usepackage{xcolor}
\usepackage{tikz}
\usetikzlibrary{arrows.meta,positioning,shapes.geometric,calc,fit,backgrounds}

%------------------------------- Tablas --------------------------------
\usepackage{booktabs}
\usepackage{tabularx}
\usepackage{array}
\usepackage{longtable}
\usepackage{multirow}
\newcolumntype{Y}{>{\raggedright\arraybackslash}X}
\newcolumntype{L}[1]{>{\raggedright\arraybackslash}p{#1}}

%------------------------------- Listas --------------------------------
\usepackage{enumitem}
\setlist{nosep,leftmargin=1.4em}

%---------------------------- Encabezados ------------------------------
\usepackage{fancyhdr}
\usepackage{titlesec}
\pagestyle{fancy}
\fancyhf{}
\fancyhead[L]{\small\leftmark}
\fancyfoot[C]{\thepage}
\renewcommand{\headrulewidth}{0.4pt}

%--------------------- Caja para TRABAJO PENDIENTE ---------------------
\usepackage[most]{tcolorbox}
\newtcolorbox{pendientebox}{
  colback=yellow!12,
  colframe=orange!75!black,
  boxrule=0.6pt,
  arc=2pt,
  left=6pt,right=6pt,top=4pt,bottom=4pt,
  fonttitle=\bfseries,
  title={\faClipboard\ Sección pendiente — trabajo futuro}
}
% comando corto: \pendiente{texto con lo que falta}
\newcommand{\pendiente}[1]{%
  \begin{pendientebox}#1\end{pendientebox}}
% si no se desea cargar fontawesome, se define un símbolo neutro:
\providecommand{\faClipboard}{$\square$}

%--------------------- Caja para NOTA DE VERSIÓN ----------------------
\newtcolorbox{notaversionbox}{
  colback=blue!6,
  colframe=blue!55!black,
  boxrule=0.6pt,
  arc=2pt,
  left=6pt,right=6pt,top=4pt,bottom=4pt,
  fonttitle=\bfseries,
  title={Estado del proceso}
}
\newcommand{\notaversion}[1]{%
  \begin{notaversionbox}#1\end{notaversionbox}}

%------------------------------ Hipervínculos --------------------------
\usepackage[hidelinks]{hyperref}
\hypersetup{
  pdftitle={THUNDERS-Master},
  pdfauthor={Andres Felipe Collazos Fernandez; Jose David Chilito Cometa},
  bookmarksnumbered=true
}

%------------------- Numeración de actividades A1..A11 -----------------
\newcommand{\act}[1]{\textbf{#1}}

\begin{document}

%=======================================================================
%  PORTADA
%=======================================================================
\begin{titlepage}
\centering
{\scshape\large Universidad del Cauca\par}
{\scshape Facultad de Ingeniería Electrónica y Telecomunicaciones\par}
{\scshape Programa de Ingeniería de Sistemas\par}
\vspace{1.5cm}

\rule{\linewidth}{0.4pt}\\[0.4cm]
{\LARGE\bfseries THUNDERS-Master:\\[0.3cm]
un proceso ligero para construir, monitorear y validar el entendimiento
compartido en las actividades de elicitación, especificación y validación
de requisitos de software\par}
\rule{\linewidth}{0.4pt}\\[1.2cm]

{\large Trabajo de grado presentado como requisito parcial\\
para optar al título de Ingeniero de Sistemas\par}
\vspace{1.5cm}

{\large\bfseries Autores\par}
{\large Andrés Felipe Collazos Fernández\par}
{\large José David Chilito Cometa\par}
\vspace{1.0cm}

{\large\bfseries Director\par}
{\large Vanessa Agredo-Delgado, Ph.D.\par}
\vspace{0.3cm}
{\large\bfseries Codirector\par}
{\large Cesar Alberto Collazos, Ph.D.\par}
\vspace{0.3cm}
{\large\bfseries Codirector\par}
{\large Leandro Antonelli, Ph.D.\par}

{\large Línea de investigación: Ingeniería de Software\par}
\vspace{0.6cm}
{\large Popayán, Colombia\par}
{\large \the\year\par}
\end{titlepage}

\pagenumbering{roman}

%=======================================================================
%  DEDICATORIA  (pendiente)
%=======================================================================
\chapter*{Dedicatoria}
\addcontentsline{toc}{chapter}{Dedicatoria}
\thispagestyle{plain}
\pendiente{Escribir la dedicatoria personal de los autores.}

%=======================================================================
%  AGRADECIMIENTOS  (pendiente)
%=======================================================================
\chapter*{Agradecimientos}
\addcontentsline{toc}{chapter}{Agradecimientos}
\thispagestyle{plain}
\pendiente{Redactar los agradecimientos (familia, director, codirector, grupo IDIS,
universidad, colaboradores y demás personas o instituciones que apoyaron el trabajo).}

%=======================================================================
%  RESUMEN
%=======================================================================
\chapter*{Resumen}
\addcontentsline{toc}{chapter}{Resumen}
\thispagestyle{plain}

El entendimiento compartido es una condición crítica para la calidad de los
requisitos de software, porque permite que stakeholders, usuarios, analistas
y equipos técnicos converjan en una interpretación común del problema, las
necesidades, los requisitos y los criterios de validación. Sin embargo, la
literatura aborda este fenómeno de forma dispersa, bajo términos como
comunicación, coordinación, alineación, negociación, entendimiento común y
modelos mentales compartidos, lo que dificulta usarlo como guía operativa
para la elicitación, especificación y validación de requisitos. Este trabajo
de grado sintetiza una revisión sistemática de la literatura y una
caracterización de dichas actividades con base en la norma
ISO/IEC/IEEE~29148:2018, y a partir de allí adapta el proceso THUNDERS hacia
\emph{THUNDERS-Master}, una versión ligera orientada a construir, monitorear
y validar el entendimiento compartido en ingeniería de requisitos. La
adaptación se realiza mediante Ingeniería de Métodos Situacionales (SME)
orientada por \emph{method chunks}, complementada con criterios de
\emph{tailoring} contextual y reducción de carga cognitiva. THUNDERS-Master
conserva la lógica de tres fases de THUNDERS —Beginning, Developing y
Measuring— pero la traduce a un proceso mínimo viable de once actividades
operativas (A1--A11), con roles simplificados, artefactos mínimos, evidencias
observables, criterios de activación contextual y trazabilidad explícita con
el proceso original. El resultado se presenta como una especificación formal
de trabajo, lista para revisión experta, walkthrough metodológico y piloto
posterior; aún no se presenta como un proceso empíricamente validado.

\vspace{0.4cm}
\noindent\textbf{Palabras clave:} Ingeniería de requisitos, Entendimiento
compartido, Ingeniería de Métodos Situacionales, Adaptación de procesos,
THUNDERS, THUNDERS-Master.

%=======================================================================
%  ABSTRACT
%=======================================================================
\chapter*{Abstract}
\addcontentsline{toc}{chapter}{Abstract}
\thispagestyle{plain}

Shared understanding is a critical condition for software requirements
quality, because it allows stakeholders, users, analysts, and technical teams
to converge on a common interpretation of the problem, the needs, the
requirements, and the validation criteria. However, the literature addresses
this phenomenon in a dispersed way, under terms such as communication,
coordination, alignment, negotiation, common understanding, and shared mental
models, which makes it difficult to use as operational guidance for
requirements elicitation, specification, and validation. This undergraduate
thesis synthesizes a systematic literature review and a characterization of
these activities based on the ISO/IEC/IEEE~29148:2018 standard, and from there
adapts the THUNDERS process into \emph{THUNDERS-Master}, a lightweight version
aimed at building, monitoring, and validating shared understanding in
requirements engineering. The adaptation is carried out through Situational
Method Engineering (SME) guided by method chunks, complemented with contextual
tailoring criteria and cognitive-load reduction. THUNDERS-Master preserves the
three-phase logic of THUNDERS —Beginning, Developing, and Measuring— but
translates it into a minimum viable process of eleven operational activities
(A1--A11), with simplified roles, minimum artifacts, observable evidence,
contextual activation criteria, and explicit traceability to the original
process. The result is presented as a formal working specification, ready for
expert review, methodological walkthrough, and subsequent piloting; it is not
yet presented as an empirically validated process.

\vspace{0.4cm}
\noindent\textbf{Keywords:} Requirements engineering, Shared understanding,
Situational Method Engineering, Process adaptation, THUNDERS, THUNDERS-Master.

%=======================================================================
%  ÍNDICES
%=======================================================================
\tableofcontents
\listoffigures
\listoftables

\cleardoublepage
\pagenumbering{arabic}

%=======================================================================
%  CAPÍTULO 1 — INTRODUCCIÓN
%=======================================================================
\chapter{Introducción}
\label{ch:introduccion}

\section{Generalidades}
La ingeniería de requisitos (IR) cumple una función mediadora entre el dominio
del problema y el dominio de la solución. Su propósito es descubrir, negociar,
documentar, verificar, validar y gestionar requisitos que expresen necesidades
reales de los stakeholders, restricciones del contexto y condiciones de
aceptación establecidas por estándares internacionales~\cite{iso29148},
actuando como una hoja de ruta para todo el ciclo de
desarrollo~\cite{nuseibeh2000roadmap}. Esta función no es únicamente técnica:
en una sesión de requisitos convergen usuarios finales, clientes,
representantes de negocio, analistas, desarrolladores, testers, operadores y
expertos de dominio, cada uno con vocabularios, prioridades y niveles de
autoridad diferentes. Por ello, la calidad de los requisitos depende tanto de
la forma del artefacto como del grado en que los participantes logran compartir
el significado que dicho artefacto representa~\cite{fuller2021blurring}.

En este trabajo, el \emph{entendimiento compartido} se entiende como el grado
en que los participantes convergen en la interpretación de un objeto de
entendimiento, comparten una perspectiva común y coordinan acciones alrededor
de un objetivo común~\cite{bittner2014shared}. En IR, ese objeto puede ser un
problema, un concepto del dominio, una necesidad, un requisito, una
restricción, un supuesto, un escenario, una historia de usuario o un criterio
de aceptación. Cuando esa convergencia es insuficiente, el equipo puede creer
que está de acuerdo aunque cada participante interprete de manera distinta el
alcance, el valor, la prioridad o la forma de validar un requisito.

\section{Motivación}
Las fallas de entendimiento compartido no siempre son visibles de inmediato.
Pueden quedar ocultas tras una ambigüedad persistente en las
especificaciones~\cite{asadabadi2020ambiguity}, supuestos no declarados,
conflictos no resueltos, pérdida de contexto, requisitos incompletos o
problemas operativos derivados de una trazabilidad
débil~\cite{fucci2022traceability}. En la elicitación, estas fallas aparecen
cuando se consulta a los stakeholders incorrectos o cuando las necesidades se
capturan como frases aisladas sin fuente, motivación ni contexto. En la
especificación, se manifiestan cuando las necesidades se transforman en
requisitos ambiguos, no verificables o desconectados de su origen. En la
validación, surgen cuando los participantes aprueban un requisito por su forma
textual, pero no porque todos lo interpreten de la misma manera.

THUNDERS ofrece una base conceptual relevante para atender este problema,
porque fue diseñado para mejorar actividades colaborativas mediante la
construcción, el monitoreo y la asistencia del entendimiento
compartido~\cite{agredo2022thesis}. Su estructura considera fases, tareas,
roles, productos de trabajo, mecanismos de monitoreo y guías de asistencia.
Sin embargo, su aplicación directa a la IR no es trivial: las validaciones
previas muestran que el proceso original es completo y metodológicamente
robusto, pero en entornos dinámicos resulta extenso, difícil de aplicar y
cognitivamente demandante~\cite{agredo2022exploratory}. En contextos de
requisitos —donde muchas actividades ocurren en talleres breves, reuniones de
refinamiento, entrevistas o revisiones con stakeholders— un proceso pesado
dificulta la adopción.

\section{Planteamiento del problema}

\subsection{Descripción del problema}
La ingeniería de requisitos es una función interdisciplinaria que media entre
los dominios del adquirente y el proveedor para establecer y mantener los
requisitos que debe cumplir un sistema, software o servicio de
interés~\cite{iso29148}. Esta función comprende actividades sistemáticas de
descubrimiento, elicitación, análisis, especificación, verificación,
validación, comunicación, documentación y gestión de requisitos, que se
ejecutan de forma iterativa a lo largo del ciclo de vida y que demandan la
participación coordinada de clientes, usuarios, desarrolladores, operadores y
responsables de negocio para construir y mantener un entendimiento compartido
de las necesidades y restricciones del sistema. Sin este entendimiento, los
grupos suelen tener interpretaciones distintas de lo que deben hacer y del
problema a resolver, lo que genera baja colaboración y resultados no
deseados~\cite{agredo2022thesis}.

Construir este entendimiento es difícil debido a la diversidad de perspectivas,
conocimientos y experiencias de quienes definen los requisitos, a la
complejidad de los contextos y conceptos que cada actor maneja, a la carga
cognitiva asociada al esfuerzo de interpretar, coordinar y mantener una
comprensión común, y a la ausencia de procesos que aseguren su mantenimiento en
el tiempo~\cite{agredo2022thesis}. Esta dificultad no surge de forma aislada:
se ve influida por condiciones organizacionales, culturales y técnicas. Durante
la elicitación y el análisis pueden aparecer criterios con definiciones
ambiguas o contradictorias que dificultan acuerdos estables~\cite{mattmann2015quality};
factores culturales como la \emph{distancia de poder} limitan la participación
de los stakeholders con menor jerarquía y sesgan la negociación~\cite{alsanoosy2019power};
y, en equipos distribuidos, las distancias geográfica, temporal y sociocultural
exigen mecanismos explícitos para sostener una comprensión común del producto,
sus requisitos y su entorno~\cite{grande2024devops}.

Frente a esta dificultad se retoma el proceso THUNDERS, construido mediante
Ingeniería de Métodos Situacionales y especificado en SPEM~2.0 a dos niveles
—conceptual y tecnológico— para diseñar, ejecutar y validar actividades de
resolución de problemas con grupos heterogéneos~\cite{agredo2022thesis}. Las
experiencias de validación reportadas muestran que THUNDERS es completo y útil
y que apoya la construcción de entendimiento compartido; no obstante, también
es un proceso largo, difícil de usar y generador de alta carga cognitiva, que
aún requiere identificar las tareas obligatorias según el contexto y derivar
versiones más simplificadas que conserven sus aportes pero con una complejidad
adecuada para entornos reales~\cite{agredo2022thesis,agredo2022exploratory}.
El problema, por tanto, no es la ausencia de mecanismos en la literatura, sino
la falta de una ruta operativa, ligera y trazable que permita construir,
monitorear y validar el entendimiento compartido durante la elicitación, la
especificación y la validación de requisitos.

\subsection{Pregunta de investigación}
\emph{¿Cómo puede apoyarse la construcción y el mantenimiento del entendimiento
compartido durante las actividades de elicitación, especificación y validación
de requisitos de un producto software, a partir del proceso THUNDERS?}

\section{Justificación}

\subsection{¿Por qué se necesita un proceso para el entendimiento compartido en IR?}
En la práctica, es frecuente que el entendimiento compartido se pierda o resulte
incompleto, porque los actores tienen diferencias de conocimiento, experiencia,
lenguaje técnico y contexto~\cite{grande2024devops}. Estas diferencias llevan a
interpretaciones distintas de lo que se quiere lograr y, ante la ausencia de un
seguimiento sistemático de lo entendido y acordado, el grupo avanza con
concepciones parciales o contradictorias que derivan en retrabajo. No basta con
documentar los requisitos: es necesario asegurar que todas las partes comprendan
lo mismo. Sin un proceso explícito que promueva y verifique esta alineación, el
riesgo de malentendidos persiste a lo largo de toda la IR~\cite{iso29148}.

\subsection{¿Por qué elicitación, especificación y validación?}
Estas tres actividades constituyen el núcleo donde se construye, negocia y
consolida el significado de los requisitos. En la elicitación los actores
expresan necesidades inicialmente implícitas, incompletas o contradictorias; en
la especificación esas necesidades se transforman en representaciones que buscan
reducir ambigüedad; y en la validación se confirma que lo documentado refleja la
intención de los interesados y que todos interpretan los requisitos de manera
coherente. Actividades posteriores como diseño, implementación y pruebas
materializan y verifican una solución basada en requisitos ya definidos, pero no
son el espacio óptimo para construir entendimiento compartido desde cero:
detectar discrepancias allí implica mayores costos de corrección y
retrabajo~\cite{agredo2022thesis,iso29148}.

\subsection{¿Por qué usar THUNDERS?}
THUNDERS ofrece una base compuesta por roles, tareas, productos de trabajo,
momentos de medición y mecanismos de validación, cuyos resultados han mostrado
mejoras en la colaboración y en el logro del entendimiento
compartido~\cite{agredo2022thesis}. Esto lo convierte en un punto de partida
sólido en lugar de diseñar un proceso completamente nuevo. La justificación no
es crear un proceso desde cero, sino adaptar y simplificar los aportes de
THUNDERS para ajustarlos a la IR en contextos reales, garantizando que el
entendimiento compartido sea visible, comprobable y sostenible.

\subsection{¿Por qué validar en un contexto agrícola?}
\label{sec:contexto_agricola}
Se propone validar el proceso en un contexto agrícola por las particularidades
comunicativas, colaborativas y sociotécnicas del sector agropecuario, donde la
adopción de nuevas herramientas depende de la capacidad de los actores para
construir significados compartidos y consensuar decisiones sobre los requisitos.
Frente a los desafíos del cambio climático y la incorporación de tecnologías
digitales, estos entornos exigen interacciones efectivas entre productores,
asesores, técnicos, investigadores y diseñadores, con diferencias de lenguaje,
marcos conceptuales y percepciones del riesgo que incrementan la probabilidad de
malentendidos~\cite{getson2022communication}. Se ha documentado que múltiples
proyectos tecnológicos del sector fracasan por no considerar adecuadamente las
percepciones, lenguajes y contextos de los usuarios finales~\cite{lopez2021crowdre},
y que una comprensión colectiva explícita favorece la toma de decisiones, la
confianza y la adopción de tecnologías resilientes~\cite{sulaiman2021ict}. Por
ello, este dominio fortalece la utilidad, transferibilidad y adaptabilidad del
proceso propuesto.

\section{Objetivos}

\subsection{Objetivo general}
Adaptar el proceso THUNDERS para apoyar la construcción del entendimiento
compartido en las actividades de elicitación, especificación y validación de
requisitos de la ingeniería de requisitos, integrándolo como apoyo al proceso
de desarrollo de un producto de software.

\subsection{Objetivos específicos}
\begin{enumerate}
  \item Caracterizar, de acuerdo con la literatura, los elementos de las
        actividades de elicitación, especificación y validación de requisitos
        que requieren construir entendimiento compartido.
  \item Identificar del proceso THUNDERS los componentes, principios y
        mecanismos susceptibles de adaptación para apoyar las actividades de
        elicitación, especificación y validación de la IR.
  \item Definir THUNDERS-Master como un proceso adaptado que permita construir
        el entendimiento compartido durante las actividades de elicitación,
        especificación y validación de requisitos de software.
  \item Evaluar, a través de un estudio de caso en un contexto agrícola, la
        utilidad del proceso adaptado para la construcción del entendimiento
        compartido en un equipo de desarrollo de un producto de software
        orientado a dicho contexto.
\end{enumerate}

\section{Metodología de la investigación}
\label{sec:metodologia}
La investigación adopta una metodología de carácter aplicado, iterativo y
orientado a la evidencia~\cite{petersen2008systematic}, adecuada para adaptar y
validar procesos en contextos reales. La iteración se materializa en la
posibilidad de refinar la propuesta con base en la evidencia recogida durante el
pilotaje. Para definir y justificar las métricas de evaluación se adopta el
enfoque Goal-Question-Metric (GQM)~\cite{basili1994gqm}, que permite
operacionalizar el constructo de entendimiento compartido en indicadores
observables y medibles, garantizando trazabilidad entre objetivos, preguntas de
evaluación y métricas. El trabajo se organiza en cuatro fases.

\begin{description}
  \item[Fase 1 — Exploración.] Revisión sistemática de la literatura (RSL) para
        caracterizar los elementos de elicitación, especificación y validación
        que demandan entendimiento compartido (Objetivo~1), y análisis documental
        de THUNDERS para identificar los componentes adaptables a IR
        (Objetivo~2). Su producto es un informe que consolida la caracterización
        y el mapeo de elementos de THUNDERS.
  \item[Fase 2 — Formulación del proceso.] Búsqueda dirigida de modelos de
        adaptación, selección del más pertinente y su aplicación estructurada
        (criterios de adaptación, mapeo de elementos, reglas de transformación y
        especificación de la versión adaptada), junto con el diseño del plan de
        medición GQM (Objetivo~3). La fase cierra con un taller de validación con
        expertos y usuarios clave para refinar la versión a pilotar.
  \item[Fase 3 — Evaluación.] Implementación piloto de THUNDERS-Master en un
        estudio de caso en contexto agrícola (Objetivo~4), mediante sesiones de
        elicitación, especificación y validación, ejecutando el plan GQM. Los
        datos se analizan con estadística descriptiva (percepción) y análisis
        temático (información cualitativa), produciendo un registro de lecciones
        aprendidas y ajustes al proceso.
  \item[Fase 4 — Documentación y transferencia.] Consolidación de los resultados
        en la monografía y sustentación del trabajo de grado.
\end{description}

\section{Aportes del proyecto}
El aporte se centra en definir un proceso adaptado para construir entendimiento
compartido específicamente en elicitación, especificación y validación de
requisitos. THUNDERS-Master no sustituye prácticas existentes: las complementa
con principios de ingeniería de la colaboración y análisis contextual. En el
ámbito académico, amplía la comprensión teórica y práctica del entendimiento
compartido en IR; en el industrial, aporta un proceso formal pero liviano que
reduce ambigüedades, retrabajo y errores de interpretación; y en el
investigativo, extiende la aplicación de THUNDERS al dominio de la IR, generando
evidencia empírica sobre cómo el entendimiento compartido puede construirse,
monitorearse y evaluarse. La validación en un contexto agrícola aporta una
perspectiva aplicada en un escenario con alta diversidad técnica, disciplinar y
comunicativa~\cite{sulaiman2021ict,lopez2021crowdre}.

\section{Organización del documento}
El resto del documento se organiza así. El Capítulo~\ref{ch:marco} presenta el
marco teórico, los trabajos relacionados y la revisión sistemática de la
literatura. El Capítulo~\ref{ch:caracterizacion} caracteriza las actividades de
elicitación, especificación y validación. El Capítulo~\ref{ch:adaptacion}
describe la Ingeniería de Métodos Situacionales y las decisiones de adaptación
de THUNDERS. El Capítulo~\ref{ch:proceso} especifica THUNDERS-Master como
proceso operativo (fases, actividades, roles, artefactos y variabilidad
contextual). El Capítulo~\ref{ch:evaluacion} presenta el plan y los resultados
de la evaluación. Finalmente, el Capítulo~\ref{ch:conclusiones} expone
conclusiones, contribuciones, limitaciones y trabajo futuro.

%=======================================================================
%  CAPÍTULO 2 — MARCO CONCEPTUAL Y REVISIÓN SISTEMÁTICA
%=======================================================================
\chapter{Marco conceptual y estado del arte}
\label{ch:marco}

\section{Generalidades}
Este capítulo establece las bases conceptuales del trabajo y sintetiza la
evidencia disponible. Primero define los conceptos de ingeniería de requisitos
y entendimiento compartido; luego describe THUNDERS como proceso base; y
finalmente resume la revisión sistemática de la literatura (RSL) que sirvió
como insumo de diseño para THUNDERS-Master.

\section{Marco teórico}

\subsection{Ingeniería de requisitos}
La ingeniería de requisitos (IR) es un proceso sistemático y repetible para
identificar, analizar, documentar y validar lo que un sistema debe hacer y bajo
qué condiciones, de modo que todos los interesados compartan una base común para
el desarrollo y la aceptación del sistema. La norma ISO/IEC/IEEE~29148:2018
define la IR como el conjunto de actividades que transforman las necesidades de
los interesados en requisitos verificables y gestionables a lo largo del ciclo
de vida —elicitar, analizar/negociar, especificar, validar y
gestionar—~\cite{iso29148}. Este trabajo se delimita a elicitación,
especificación y validación porque son los momentos donde el significado de los
requisitos se construye, se formaliza y se confirma.

\subsection{Elicitación}
Es la actividad en la que se identifican y recopilan las necesidades,
expectativas y restricciones de los stakeholders sobre el sistema. Según
ISO/IEC/IEEE~29148, busca capturar requisitos desde distintas fuentes y
perspectivas, por lo que requiere preparación e identificación de fuentes para
obtener información que sirva de base para su análisis y
documentación~\cite{iso29148}. No es una simple recolección de información, sino
una construcción inicial de significado~\cite{lim2021data}.

\subsection{Especificación}
Corresponde a la actividad en la que los requisitos elicitados se documentan y
estructuran de forma clara, completa y verificable, para servir como base de
comunicación y acuerdo entre stakeholders y equipo de desarrollo. Los requisitos
se formulan con redacción precisa y medible, manteniendo consistencia y
trazabilidad para facilitar su revisión y uso en etapas
posteriores~\cite{iso29148,behutiye2022quality}.

\subsection{Validación}
Es la actividad mediante la cual se confirma, junto con los stakeholders, que los
requisitos documentados representan correctamente sus necesidades y son adecuados
para el propósito del sistema. Se realiza a través de revisiones y
retroalimentación con los interesados, con el fin de detectar inconsistencias,
ambigüedades u omisiones antes de avanzar~\cite{iso29148}.

\subsection{Trabajo colaborativo}
El trabajo colaborativo se entiende como la contribución conjunta de ideas,
experiencias y conocimientos de un grupo orientado a una meta común. Su propósito
no es solo optimizar resultados, sino generar conocimiento colectivo y fortalecer
el entendimiento compartido del problema~\cite{agredo2022thesis}. Se caracteriza
por la interdependencia positiva: el éxito depende de integrar los aportes en una
construcción conjunta, más allá del reparto de tareas, y requiere comunicación
efectiva, escucha activa y toma de decisiones compartida~\cite{spijkman2021alignment}.

\subsection{Entendimiento compartido}
El entendimiento compartido (\emph{Shared Understanding}, SU) se construye a
partir del intercambio de información y la interacción social, y se refleja en
acuerdos sobre procesos, significados y orden de las actividades, permitiendo
coordinar conductas hacia objetivos comunes~\cite{bittner2014shared,agredo2022thesis}.
En IR se materializa cuando clientes, usuarios, analistas, diseñadores y demás
interesados interpretan de manera equivalente las necesidades, restricciones y
criterios de calidad, lo que reduce ambigüedades y retrabajo~\cite{mattmann2015quality}.
La literatura lo aborda de forma dispersa bajo términos como comunicación,
coordinación, alineación, negociación, entendimiento común y modelos mentales
compartidos, lo que limita su uso como guía operativa.

\subsection{Goal-Question-Metric (GQM)}
GQM es un método dirigido por objetivos para definir e interpretar mediciones en
ingeniería de software: primero se establece un objetivo de medición, luego se
descompone en preguntas que permiten evaluar su cumplimiento y, finalmente, se
definen métricas concretas para responder esas preguntas. Así, asegura que las
métricas recolectadas estén alineadas con las necesidades de información del
proyecto~\cite{basili1994gqm}.

\subsection{THUNDERS como proceso base}
THUNDERS es un proceso de trabajo colaborativo orientado a construir, medir,
monitorear y asistir el entendimiento compartido durante actividades de
resolución de problemas~\cite{agredo2022thesis}. Fue construido con Ingeniería de
Métodos Situacionales y formalizado en SPEM~2.0, a nivel conceptual y
tecnológico, y organiza fases, actividades, roles, tareas, productos de trabajo y
mecanismos de validación. Su estructura se organiza en tres fases:
\emph{Beginning} (diseñar y especificar la actividad colaborativa),
\emph{Developing} (ejecutar la actividad y construir el entendimiento) y
\emph{Measuring} (validar la solución, evaluar desempeño y cerrar), con
mecanismos transversales de monitoreo y asistencia. No obstante, su validación
reporta que, aunque es completo y útil, también es extenso, difícil de aplicar y
generador de alta carga cognitiva~\cite{agredo2022exploratory}.

\subsection{Adaptación de procesos}
La adaptación de procesos es el ajuste sistemático de un proceso para adecuarlo a
las condiciones específicas de un proyecto u organización, manteniendo su
eficacia ante variables como recursos, cultura, capacidades del equipo,
herramientas o requisitos de calidad~\cite{li2024adaptability}. Adaptar no es
solo modificar tareas o roles, sino asegurar consistencia con los objetivos del
proceso. Entre los pasos comúnmente considerados se incluyen: diagnóstico y
contextualización, definición de criterios de adaptación, mapeo del proceso
original, diseño de reglas de transformación, documentación del proceso adaptado
y validación con mejora iterativa~\cite{li2024adaptability}.

\section{Trabajos relacionados}
La revisión de la literatura identificó investigaciones que, desde distintos
enfoques, buscan construir, mantener y medir el entendimiento compartido en la IR.

\subsection{Ingeniería de requisitos y entendimiento compartido}
Los trabajos de esta línea muestran que contar con artefactos compartidos y
representaciones comunes mejora la comunicación y reduce la ambigüedad en las
fases iniciales. El uso de \emph{personas} y escenarios narrativos facilita
entender las motivaciones de los usuarios~\cite{wang2025personas}, y acordar y
mantener visibles los criterios de calidad disminuye errores de interpretación y
deuda técnica~\cite{behutiye2022quality}. Sin embargo, muchos enfoques abordan el
SU de manera fragmentada —centrados en artefactos, modelos o dinámicas sociales
por separado— sin integrarlos en un proceso coherente y continuo.

\subsection{Medición del entendimiento en el ciclo de vida del desarrollo}
Una segunda línea estudia cómo evidenciar o medir el SU según la fase del ciclo
de vida. En fases tempranas se emplean grafos de conocimiento y modelos
semánticos para evaluar si los actores comparten significados equivalentes sobre
el dominio; al avanzar, la medición se apoya en la trazabilidad y la consistencia
entre artefactos. Pese a ello, la literatura coincide en que aún no existe una
metodología estandarizada para un seguimiento continuo del SU dentro de la
IR~\cite{agredo2022thesis,iso29148}, lo que justifica incorporar mecanismos
formales de monitoreo y verificación.

\subsection{Construcción del entendimiento en metodologías de desarrollo}
Los paradigmas de desarrollo influyen en cómo los equipos construyen y mantienen
el SU. Los enfoques empíricos y ágiles ofrecen más oportunidades para
co-construir modelos mentales compartidos~\cite{ralph2018paradigms}, pero los
estudios sobre adopciones ágiles a gran escala muestran que, sin estructuras
claras de coordinación, trazabilidad y roles de apoyo, el entendimiento tiende a
fragmentarse entre equipos y niveles organizacionales~\cite{dikert2016largescale}.
El reto no es oponer «ágil» frente a «tradicional», sino disponer de procesos que
hagan explícita la construcción, el monitoreo y la verificación del SU.

\subsection{Colaboración distribuida y factores culturales}
Los estudios sobre colaboración distribuida examinan cómo la dispersión
geográfica, los husos horarios y las diferencias culturales afectan la
coordinación. Cuando los miembros trabajan desde distintos lugares surgen riesgos
de pérdida de contexto, retrasos en la retroalimentación y desalineación del
SU~\cite{grande2024devops}. Prácticas de integración continua, repositorios
colaborativos y reuniones de sincronización ayudan a mitigarlos. En cuanto a lo
cultural, las diferencias jerárquicas y de poder influyen en la participación y
la calidad de la negociación: las estructuras rígidas limitan la
retroalimentación, mientras que las culturas más horizontales fomentan la
confianza y la validación colectiva~\cite{alsanoosy2019power}.

\subsection{Adaptación de procesos y experiencias contextuales}
La adaptación de procesos es una estrategia clave para ajustar metodologías a
contextos ágiles o de pequeña escala, evaluando los cambios no solo por su
impacto técnico sino por su influencia en la coordinación y el entendimiento del
equipo~\cite{dikert2016largescale,li2024adaptability}. De forma complementaria,
los resultados de THUNDERS en escenarios colaborativos son positivos, pero
evidencian la necesidad de versiones más ligeras y adaptadas a equipos reales,
con artefactos simplificados y mecanismos de validación
claros~\cite{agredo2022thesis}. Esta brecha es justamente la que THUNDERS-Master
busca atender.

\section{Revisión sistemática de la literatura}
\label{sec:rsl}
La RSL se condujo siguiendo el proceso de mapeo sistemático propuesto por
Petersen et al.~\cite{petersen2008systematic}. Su objetivo fue identificar,
organizar y sintetizar estudios primarios que reportan rupturas, situaciones
críticas, mecanismos de apoyo y formas de evaluación del entendimiento
compartido en la elicitación, especificación y validación de requisitos.

\subsection{Protocolo y proceso de revisión}
La revisión siguió un protocolo definido a priori, organizado en las fases del
mapeo sistemático~\cite{petersen2008systematic}: \emph{(i)}~planificación, que fijó
el objetivo, las preguntas de investigación y los criterios de inclusión y
exclusión; \emph{(ii)}~búsqueda, con una cadena estructurada aplicada a las bases
seleccionadas; \emph{(iii)}~cribado en dos etapas, primero por título y resumen y
luego por texto completo; \emph{(iv)}~extracción de datos mediante un esquema
alineado con las preguntas de investigación; y \emph{(v)}~síntesis temática de los
estudios primarios. La búsqueda, los criterios y el flujo de selección se detallan
en las subsecciones siguientes; aquí se describe el cribado y la extracción que
sostienen la trazabilidad entre la evidencia y los hallazgos.

El cribado se aplicó de forma secuencial. En la primera etapa se evaluaron título,
resumen y palabras clave frente a los criterios de inclusión y exclusión, y se
descartaron los estudios claramente fuera del dominio de requisitos o sin relación
con el entendimiento compartido. En la segunda etapa se leyó el texto completo de
los estudios restantes y se conservaron solo aquellos que aportaban evidencia útil
para al menos una pregunta de investigación. Ambas etapas fueron realizadas de forma
independiente por los dos autores del trabajo; los desacuerdos sobre la inclusión de
un estudio se resolvieron por discusión directa entre ambos y, cuando el caso lo
ameritaba, mediante consulta a la directora del trabajo de grado. Adicionalmente, la
directora revisó el proceso de cribado y sus resultados como control de calidad de
la selección.
La extracción se apoyó en un esquema que asocia cada campo con la pregunta de
investigación que alimenta (Tabla~\ref{tab:rsl_extraction}), de modo que la síntesis
por pregunta de investigación se construyó sobre datos recolectados de forma
homogénea y no sobre lecturas \emph{ad hoc} de cada estudio.

\begin{table}[htbp]
\caption{Esquema de extracción de datos y su relación con las preguntas de investigación.}
\label{tab:rsl_extraction}
\centering
\footnotesize
\renewcommand{\arraystretch}{1.15}
\begin{tabularx}{\textwidth}{@{}L{4.6cm}Y L{1.7cm}@{}}
\toprule
\textbf{Campo extraído} & \textbf{Descripción} & \textbf{Pregunta} \\
\midrule
Metadatos del estudio & Autores, año, fuente, tipo de estudio y actividad(es) de IR
abordada(s). & Todas \\
Problema o ruptura de entendimiento & Fallas de comunicación, ambigüedad o
desalineación reportadas. & RQ1 \\
Tarea, artefacto o situación & Dónde se identifica la necesidad de construir
entendimiento compartido. & RQ2 \\
Momento crítico & Punto del proceso donde el significado es inestable y debe
negociarse. & RQ3 \\
Mecanismo, técnica o enfoque & Solución propuesta o usada para construir o mantener
el entendimiento. & RQ4 \\
Métrica, indicador o instrumento & Forma de medir o evaluar el entendimiento y su
momento de aplicación. & RQ5 \\
\bottomrule
\end{tabularx}
\end{table}

\subsection{Preguntas de investigación}
\begin{itemize}
  \item \textbf{RQ1.} ¿Qué problemas de comunicación o desalineaciones del
        entendimiento compartido se reportan durante la elicitación,
        especificación y validación?
  \item \textbf{RQ2.} ¿En qué tareas, artefactos y situaciones se identifica la
        necesidad de construir entendimiento compartido entre stakeholders?
  \item \textbf{RQ3.} ¿Cuáles son los momentos más críticos para lograr
        entendimiento compartido en cada actividad y por qué?
  \item \textbf{RQ4.} ¿Qué mecanismos, técnicas o enfoques se han propuesto o
        usado para apoyar la construcción y el mantenimiento del entendimiento
        compartido?
  \item \textbf{RQ5.} ¿Qué métricas, indicadores e instrumentos se han propuesto
        para medir o evaluar el entendimiento compartido, y en qué momentos
        (antes, durante o después)?
\end{itemize}

\subsection{Estrategia de búsqueda y fuentes}
La búsqueda se estructuró alrededor de tres conceptos: \emph{Requirements
Engineering}, \emph{Shared Understanding} y \emph{Software Development}. La
cadena general de búsqueda fue:

\begin{quote}\small\ttfamily
\shorthandoff{"}
("Requirements Engineering" OR "Elicitation" OR\\
\hspace*{1em}"Requirements Analysis")\\
AND ("Shared Understanding" OR "Common Understanding" OR\\
\hspace*{1em}"Cognitive Alignment")\\
AND ("Software Development")
\end{quote}

La cadena no incluyó términos específicos por subactividad, con el fin de
maximizar la sensibilidad, dado que el fenómeno aparece fragmentado bajo
términos relacionados. La búsqueda se ejecutó en \emph{ScienceDirect},
\emph{Scopus} y \emph{Web of Science}, restringida a artículos de investigación
en inglés, acceso abierto y publicados entre 2015 y 2026.

\subsection{Criterios de inclusión y exclusión}
Se incluyeron estudios primarios revisados por pares que abordaran actividades
de IR y reportaran evidencia sobre entendimiento compartido o fenómenos
equivalentes entre stakeholders, con texto completo disponible y dentro del
periodo definido. Se excluyeron estudios fuera del dominio de requisitos,
revisiones secundarias, editoriales, capítulos de libro, tesis, trabajos sin
texto completo o sin evidencia útil para las preguntas, y duplicados.

\subsection{Resultados del proceso de selección}
La aplicación de la estrategia recuperó inicialmente 1.362 registros. Tras
aplicar los filtros se redujeron a 153 estudios y, eliminados los duplicados,
quedaron 144 estudios únicos (101 de ScienceDirect, 24 de Scopus y 19 de Web of
Science). La revisión de título y resumen excluyó 95 estudios y seleccionó 49
para lectura de texto completo. Finalmente se seleccionaron \textbf{29 estudios
primarios} (14 de ScienceDirect, 7 de Scopus y 8 de Web of Science). La
Figura~\ref{fig:prisma} resume el flujo de selección.

\begin{figure}[htbp]
\centering
\begin{tikzpicture}[
  node distance=0.7cm and 0cm,
  box/.style={draw,rounded corners=2pt,align=center,minimum width=8.5cm,
              minimum height=0.9cm,fill=blue!5},
  arr/.style={-{Stealth[length=2mm]},thick}]
  \node[box] (a) {Registros recuperados: \textbf{1.362}\\(ScienceDirect, Scopus, Web of Science)};
  \node[box,below=of a] (b) {Tras filtros de la estrategia: \textbf{153}};
  \node[box,below=of b] (c) {Estudios únicos tras eliminar duplicados: \textbf{144}};
  \node[box,below=of c] (d) {Seleccionados para texto completo: \textbf{49}};
  \node[box,below=of d,fill=green!10] (e) {Estudios primarios incluidos: \textbf{29}};
  \draw[arr] (a)--(b);
  \draw[arr] (b)--(c);
  \draw[arr] (c)--(d);
  \draw[arr] (d)--(e);
\end{tikzpicture}
\caption{Flujo de selección de estudios de la revisión sistemática.}
\label{fig:prisma}
\end{figure}

\subsection{Resultados por pregunta de investigación}

\subsubsection{RQ1 — Problemas de comunicación y desalineación}
Los problemas de comunicación aparecen desde las primeras interacciones y
persisten durante las tres actividades. En elicitación, el entendimiento se
debilita cuando no se identifican correctamente los stakeholders relevantes o
cuando los requisitos provienen de actores que no representan las necesidades
reales~\cite{khan2022stakeholders}, y se agrava por barreras de lenguaje,
diferencias culturales y distancia geográfica/temporal en contextos
distribuidos~\cite{grande2024devops}. En culturas con alta distancia de poder,
algunos actores evitan cuestionar decisiones de niveles
superiores~\cite{alsanoosy2019power}. En especificación, los problemas más
recurrentes son la ambigüedad, la incompletitud y la baja calidad de los
artefactos: una especificación en lenguaje natural admite múltiples
interpretaciones, de modo que autor y lectores creen compartir el mismo
significado cuando en realidad entienden cosas distintas~\cite{asadabadi2020ambiguity}.
En validación, la desalineación surge cuando lo documentado no puede contrastarse
con lo que los stakeholders quisieron expresar, por ausencia de criterios de
aceptación claros o por requisitos no verificables~\cite{fucci2022traceability}.
En suma, los problemas de comunicación no son solo fallas de redacción, sino
rupturas en la construcción, negociación y mantenimiento del significado
compartido.

\subsubsection{RQ2 — Tareas, artefactos y situaciones}
La necesidad de construir entendimiento compartido es transversal. En elicitación
aparece en la identificación de stakeholders, la recolección de necesidades, la
priorización y la clarificación de expectativas, apoyada en entrevistas,
\emph{workshops}, \emph{focus groups}, listas de stakeholders, glosarios de
dominio, prototipos y escenarios~\cite{lim2021data}. En especificación surge al
convertir lo elicitado en artefactos estructurados —historias de usuario,
escenarios \emph{Given/When/Then} de BDD, criterios de aceptación, plantillas y
modelos—~\cite{ralph2018paradigms}. En validación se identifica en revisiones
conjuntas, aprobación formal, \emph{sprint reviews} y verificación de consistencia
entre requisitos y pruebas~\cite{mattmann2015quality}. No obstante, la evidencia
muestra que estos artefactos son \emph{soportes} para externalizar y revisar el
significado, no prueba concluyente de que el entendimiento se haya logrado.

\subsubsection{RQ3 — Momentos críticos}
Los momentos más críticos corresponden a puntos donde el significado aún es
inestable y debe negociarse: la selección inicial de participantes, las sesiones
de interacción en tiempo real, la definición del vocabulario del dominio, y los
\emph{handoffs} entre quienes elicitan y quienes especifican~\cite{khan2022stakeholders}.
En especificación, el momento crítico es la transformación de necesidades en
artefactos, donde el significado deja de apoyarse en la conversación y queda
codificado para que otros lo interpreten~\cite{asadabadi2020ambiguity}. En
validación, la criticidad aparece al confrontar lo especificado con criterios de
aceptación, pruebas e implementación, y en la transferencia a actores que no
participaron en la discusión original~\cite{fucci2022traceability}. La criticidad
no depende solo del contenido, sino del momento en que el significado se aclara,
negocia, documenta o valida.

\subsubsection{RQ4 — Mecanismos, técnicas y enfoques}
La literatura reporta mecanismos sociales (entrevistas, talleres, discusiones
guiadas), técnicos (plantillas, modelos, glosarios como el LEL, trazabilidad,
\emph{backlogs}, repositorios versionados) y sociotécnicos que combinan
conversación con artefactos visibles y revisables~\cite{lim2021data,spijkman2021alignment}.
En elicitación predominan mecanismos de construcción inicial (moderador,
\emph{scribe}, pizarra compartida, \emph{grounding}, \emph{personas},
prototipos); en especificación, mecanismos de formalización (plantillas, BDD,
modelado participativo, control de ambigüedad); y en validación, mecanismos de
verificación y reajuste (criterios de aceptación, \emph{sign-off},
trazabilidad, demos, \emph{sprint reviews})~\cite{ralph2018paradigms,fucci2022traceability}.
Los mecanismos más robustos combinan interacción entre stakeholders, artefactos
estructurados, trazabilidad y retroalimentación continua, pero la literatura no
ofrece un procedimiento integrado, paso a paso, que indique cómo seleccionarlos y
combinarlos para garantizar el entendimiento compartido.

\subsubsection{RQ5 — Métricas, indicadores e instrumentos}
No existe una métrica única, directa y ampliamente aceptada para medir el
entendimiento compartido como constructo. En su lugar, se reportan indicadores
indirectos y observables distribuidos en tres momentos: \emph{pre} (líneas base,
cobertura y participación de stakeholders), \emph{durante} (solicitudes de
aclaración, malentendidos, balance de participación, señales de consenso) y
\emph{post} (cuestionarios de satisfacción y acuerdo, \emph{sign-off}, criterios
de aceptación, trazabilidad necesidad--requisito--prueba, defectos y retrabajo
posteriores)~\cite{basili1994gqm}. Estos indicadores aproximan el fenómeno, pero
no siempre distinguen entre la existencia de documentación, la aceptación formal
y el entendimiento compartido real, lo que refuerza la necesidad de instrumentos
que evalúen si el significado fue efectivamente comprendido, preservado y
validado.

\subsection{Síntesis de hallazgos}
La síntesis temática organizada por preguntas de investigación arrojó cinco
hallazgos principales que justifican las decisiones de diseño de
THUNDERS-Master. Primero, el entendimiento compartido requiere representación
adecuada de stakeholders y fuentes de conocimiento. Segundo, el lenguaje del
dominio debe gestionarse de forma explícita para evitar que términos comunes
oculten significados distintos. Tercero, las necesidades deben registrarse con
fuente, contexto, valor esperado y supuestos. Cuarto, la trazabilidad es un
mecanismo semántico —no solo administrativo— que conecta necesidad, requisito,
criterio, decisión y evidencia de validación. Quinto, la validación debe revisar
significado y criterios de aceptación, no solo completitud formal. La RSL también
muestra que no existe una definición operacional consolidada de entendimiento
compartido para IR, que pocos estudios ofrecen una guía sistemática para
construirlo y mantenerlo, y que su evaluación se apoya en indicadores indirectos
más que en métricas o instrumentos validados.

\begin{table}[htbp]
\caption{Hallazgos de la RSL y su implicación para THUNDERS-Master.}
\label{tab:rsl_implications}
\centering
\footnotesize
\renewcommand{\arraystretch}{1.15}
\begin{tabularx}{\textwidth}{@{}L{6.0cm}Y@{}}
\toprule
\textbf{Hallazgo de la RSL} & \textbf{Implicación para THUNDERS-Master} \\
\midrule
El entendimiento compartido se describe mediante conceptos dispersos
(alineación, comunicación, colaboración, modelos mentales). &
Adoptar una definición operacional y traducirla en chequeos ligeros de
alineación de interpretación. \\
\midrule
Las rupturas aparecen cuando los stakeholders relevantes no participan o su
autoridad/conocimiento no es claro. &
Incluir un mapa mínimo de stakeholders con fuente de conocimiento,
disponibilidad, representatividad y poder de decisión. \\
\midrule
La ambigüedad y los vocabularios heterogéneos afectan elicitación y
especificación. &
Mantener un glosario vivo y una actividad explícita para resolver términos
críticos o asignar responsable de aclaración. \\
\midrule
La transformación de necesidades en requisitos puede perder contexto, valor o
supuestos. &
Cada necesidad conserva fuente, contexto y estado; cada requisito crítico se
enlaza con su necesidad de origen. \\
\midrule
La validación puede volverse superficial al enfocarse solo en la forma textual. &
La validación cubre la cadena necesidad--requisito--criterio--significado y
registra las discrepancias. \\
\midrule
La literatura reporta mecanismos, pero no siempre una ruta integrada de
aplicación. &
THUNDERS-Master organiza estos mecanismos en tres fases y once actividades
operativas. \\
\bottomrule
\end{tabularx}
\end{table}

\subsection{Brecha de investigación}
\label{sec:brecha}
La evidencia sintetizada permite delimitar con precisión el problema que los
enfoques actuales no resuelven. Primero, el entendimiento compartido se aborda de
forma fragmentada: los estudios se concentran en artefactos, modelos o dinámicas
sociales por separado, sin integrarlos en una ruta continua que cubra elicitación,
especificación y validación. Segundo, no existe una definición operacional
consolidada del entendimiento compartido para IR, lo que impide tratarlo como
criterio verificable durante el trabajo. Tercero, aunque la literatura reporta
múltiples mecanismos, no ofrece un procedimiento integrado, ligero y trazable que
indique cómo seleccionarlos y combinarlos. Cuarto, su evaluación se apoya en
indicadores indirectos que no distinguen entre existencia de documentación,
aceptación formal y entendimiento real. Los marcos de desarrollo ampliamente
adoptados atienden la coordinación mediante ceremonias y refinamiento, pero no
prescriben cómo construir, monitorear y validar el significado compartido de los
requisitos ni qué evidencia mínima lo hace comprobable.

Esta brecha justifica metodológicamente la decisión de adaptar THUNDERS en lugar de
diseñar un proceso nuevo. THUNDERS fue construido con Ingeniería de Métodos
Situacionales y validado como un proceso completo y útil para construir, monitorear
y asistir el entendimiento compartido, pero también extenso y cognitivamente
demandante~\cite{agredo2022thesis,agredo2022exploratory}. Existe, por tanto, un
proceso con la intención correcta pero con una complejidad inadecuada para las
sesiones breves de la IR. La respuesta razonada no es partir de cero, sino aplicar
una adaptación situacional que conserve la intención de THUNDERS y reduzca su carga,
orientándola a las tres actividades del alcance. La Tabla~\ref{tab:gap_adaptation}
resume este razonamiento, que el Capítulo~\ref{ch:adaptacion} concreta en decisiones
de adaptación específicas.

\begin{table}[htbp]
\caption{De la brecha identificada a la justificación de adaptar THUNDERS.}
\label{tab:gap_adaptation}
\centering
\footnotesize
\renewcommand{\arraystretch}{1.15}
\begin{tabularx}{\textwidth}{@{}L{4.2cm}Y Y@{}}
\toprule
\textbf{Brecha identificada en la RSL} & \textbf{Por qué THUNDERS es una base
adecuada} & \textbf{Por qué requiere adaptación a IR} \\
\midrule
No hay una ruta integrada, ligera y trazable para construir, monitorear y validar
el entendimiento compartido. &
THUNDERS integra fases, tareas, roles, productos y mecanismos de monitoreo y
asistencia del entendimiento compartido. &
Su flujo completo es extenso para talleres, entrevistas y refinamientos; debe
simplificarse conservando la intención. \\
\midrule
Falta una definición operacional del entendimiento compartido verificable durante
el trabajo. &
THUNDERS trata el entendimiento compartido como objeto explícito de construcción y
medición. &
La operacionalización debe traducirse en chequeos ligeros y evidencia mínima
propios de la IR. \\
\midrule
Los mecanismos existen dispersos, sin criterios para seleccionarlos y combinarlos. &
THUNDERS ya organiza mecanismos sociales, técnicos y sociotécnicos en un proceso
coherente. &
La selección debe regirse por su aporte directo al entendimiento compartido en
elicitación, especificación y validación. \\
\midrule
La evaluación se apoya en indicadores indirectos que confunden documentación con
entendimiento. &
THUNDERS incorpora momentos de validación y medición del entendimiento. &
La evaluación de desempeño individual se retira del núcleo y se prioriza la
validación semántica de requisitos. \\
\bottomrule
\end{tabularx}
\end{table}

%=======================================================================
%  CAPÍTULO 3 — CARACTERIZACIÓN DE LAS ACTIVIDADES DE IR
%=======================================================================
\chapter{Caracterización de las actividades de ingeniería de requisitos}
\label{ch:caracterizacion}

\section{Generalidades}
Este capítulo caracteriza las tres actividades de IR que sirven de base para
THUNDERS-Master: elicitación, especificación y validación. La caracterización
toma como base principal la norma ISO/IEC/IEEE~29148:2018~\cite{iso29148} y se
complementa con normas de ciclo de vida, el modelo de calidad
ISO/IEC~25010~\cite{iso25010}, normas de usabilidad~\cite{iso9241} y literatura
académica. Su propósito es definir qué debe hacer el proceso, qué productos debe
generar y qué evidencias debe conservar para apoyar el entendimiento compartido.
Además de la norma, la caracterización se ancla en la evidencia de la revisión
sistemática (Sección~\ref{sec:rsl}): cada actividad se lee no solo desde su
definición normativa, sino desde las rupturas de entendimiento compartido, los
momentos críticos y los mecanismos que la RSL identificó para ella. La
Tabla~\ref{tab:re_characterization} resume la lectura operativa adoptada.

\begin{table}[htbp]
\caption{Caracterización operativa de las actividades de ingeniería de requisitos.}
\label{tab:re_characterization}
\centering
\footnotesize
\renewcommand{\arraystretch}{1.15}
\begin{tabularx}{\textwidth}{@{}L{1.9cm}L{3.0cm}Y Y@{}}
\toprule
\textbf{Actividad} & \textbf{Propósito} & \textbf{Artefactos esperados} &
\textbf{Relación con el entendimiento compartido} \\
\midrule
Elicitación &
Descubrir, aclarar, contrastar y priorizar necesidades, restricciones,
escenarios, supuestos y criterios de calidad. &
Mapa de stakeholders; contexto de uso; escenarios; registro de necesidades;
glosario inicial; conflictos y supuestos. &
Construye la primera interpretación compartida del problema y hace visibles los
desacuerdos. \\
\midrule
Especificación &
Transformar necesidades y acuerdos en requisitos bien formados, estructurados,
trazables y verificables. &
Requisitos individuales; atributos; taxonomía; matriz de evaluación; StRS, SyRS
o SRS. &
Formaliza el significado acordado, reduce la ambigüedad y preserva la
trazabilidad. \\
\midrule
Validación &
Confirmar con stakeholders que los requisitos representan el sistema correcto y
el contexto real de uso. &
Actas de revisión; comentarios; decisiones; estado de aceptación; baseline
versionada; registro de desacuerdos. &
Transforma la interpretación documentada en acuerdo explícito o desacuerdo
gestionado. \\
\bottomrule
\end{tabularx}
\end{table}

\section{Elicitación de requisitos}
La elicitación es la actividad colaborativa, iterativa y orientada al contexto
mediante la cual se identifican, contrastan, priorizan y refinan necesidades,
expectativas, restricciones, escenarios, criterios de calidad y supuestos
relevantes de los stakeholders~\cite{iso29148}. No debe confundirse con la
simple recolección de frases: su función es construir una comprensión inicial
del problema, reconocer tensiones entre actores y transformar información
dispersa en insumos útiles~\cite{nuseibeh2000roadmap}. Coherente con la RSL, aquí
se concentran las rupturas tempranas del entendimiento compartido —selección de
stakeholders no representativos y vocabulario de dominio ambiguo (RQ1, RQ3)—, por lo
que la caracterización enfatiza la identificación de fuentes y la gestión explícita
de la terminología.

\subsection{Alcance y entradas}
La elicitación cubre necesidades explícitas e implícitas, restricciones,
criterios de aceptación, escenarios de operación y atributos de calidad, e
incluye el reconocimiento de conflictos, variantes y prioridades iniciales. Debe
considerar el contexto de uso —usuarios, tareas, ambientes, herramientas,
objetivos y condiciones— y consolidar necesidades a partir de varias fuentes, sin
depender de una sola voz o técnica. Sus entradas principales son el problema u
oportunidad formulada, los objetivos de negocio, los conceptos operacionales
preliminares, la información del dominio, el inventario inicial de stakeholders,
la documentación existente y las restricciones regulatorias, técnicas y
organizacionales.

\subsection{Tareas, técnicas y productos}
Las técnicas (entrevistas semiestructuradas, talleres facilitados, co-creación,
observación y \emph{shadowing}, revisión documental, escenarios y
\emph{user journeys}, prototipos de baja fidelidad, análisis de \emph{feedback} y
datos) deben depender del tipo de fuente y de necesidad, no de una plantilla
única~\cite{lim2021data}. La Tabla~\ref{tab:tareas_elicitacion} resume las tareas
características y su evidencia esperada.

\begin{table}[htbp]
\caption{Tareas características de la elicitación.}
\label{tab:tareas_elicitacion}
\centering
\footnotesize
\renewcommand{\arraystretch}{1.12}
\begin{tabularx}{\textwidth}{@{}L{4.2cm}Y L{3.4cm}@{}}
\toprule
\textbf{Tarea} & \textbf{Descripción operativa} & \textbf{Evidencia esperada} \\
\midrule
Preparar la elicitación & Definir alcance, estrategia, fuentes, técnicas y
reglas. & Plan o agenda de elicitación. \\
Identificar stakeholders y fuentes & Reconocer usuarios, clientes, operadores,
expertos, documentos y normas. & Mapa de stakeholders y fuentes. \\
Definir contexto de uso & Describir usuarios, tareas, ambientes y modos de
operación. & Descripción de contexto y escenarios. \\
Elicitar necesidades explícitas e implícitas & Entrevistas, talleres,
observación, prototipos, análisis del dominio. & Registro de necesidades y
supuestos. \\
Analizar escenarios & Casos de uso, \emph{user journeys}, \emph{day-in-the-life}.
& Escenarios revisados con stakeholders. \\
Priorizar y depurar & Ordenar necesidades, eliminar duplicados, detectar
conflictos. & Priorización inicial y conflictos abiertos. \\
Transformar en requisitos de stakeholder & Convertir necesidades en enunciados
con fuente y contexto. & Requisitos de stakeholder preliminares. \\
Retroalimentar y negociar & Devolver la interpretación para corregir y acordar
significado. & Observaciones, acuerdos y cambios. \\
\bottomrule
\end{tabularx}
\end{table}

\subsection{Roles, artefactos y criterios de salida}
Participan el analista o facilitador, representantes de usuarios, operadores,
cliente o patrocinador, expertos de dominio y perfiles de arquitectura,
desarrollo, QA, UX o seguridad. Los artefactos incluyen mapa de stakeholders,
inventario de fuentes, descripción del contexto de uso, escenarios, registro de
necesidades, glosario inicial y registro de conflictos y supuestos. La
elicitación se considera suficientemente completa cuando hay cobertura razonable
de stakeholders, el contexto de uso está descrito, las necesidades críticas
están identificadas, los conflictos principales registrados y las necesidades
priorizadas tienen fuente y justificación.

\section{Especificación de requisitos}
La especificación transforma la base de necesidades y acuerdos en un conjunto
estructurado de requisitos bien formados, con atributos, trazabilidad y
criterios de verificación o validación~\cite{iso29148}. No es solo redacción:
implica analizar el significado de los requisitos, clasificarlos, organizarlos,
revisar su calidad y prepararlos para su uso por desarrollo, arquitectura,
pruebas y stakeholders. La RSL señala que en este punto el significado deja de
apoyarse en la conversación y queda codificado para que otros lo interpreten (RQ3),
y que la ambigüedad del lenguaje natural es la principal fuente de desalineación
(RQ1); por ello la caracterización prioriza la calidad del enunciado y la
trazabilidad hacia la necesidad de origen.

\subsection{Reglas de calidad del requisito y del conjunto}
Un requisito individual debe ser necesario, apropiado al nivel de abstracción, no
ambiguo, completo, singular, factible, verificable, correcto y conforme al estilo
adoptado. En la práctica, esto implica no mezclar varias obligaciones en una
oración, evitar palabras abiertas como «rápido» o «adecuado» sin métrica, no
incluir decisiones de diseño innecesarias y acompañar el enunciado con fuente,
justificación y método de comprobación. El conjunto debe ser completo,
consistente, factible, comprensible y validable, manteniendo trazabilidad hacia
necesidades, fuentes, decisiones y criterios de aceptación~\cite{iso29148}. Los
atributos de calidad deben tratarse de forma explícita, pues en contextos ágiles
suelen quedar incompletos o asumidos como conocimiento
tácito~\cite{behutiye2022quality}.

\subsection{Tareas, artefactos y criterios de salida}
Las tareas características son definir frontera y alcance, definir la estrategia
de especificación, formular requisitos, clasificarlos, agregar atributos
(identificador, fuente, prioridad, riesgo, \emph{owner}, \emph{rationale},
versión y método de evaluación), analizar consistencia, organizar la
documentación (StRS, SyRS, SRS o equivalente) y preparar la evaluación. Los
artefactos incluyen las especificaciones, el glosario y convenciones de
redacción, el repositorio de requisitos con atributos, la matriz de trazabilidad
y las listas de chequeo de calidad. La especificación se considera lista para
validación cuando cada requisito tiene identificador, fuente, justificación y
trazabilidad mínima; cuando los requisitos críticos tienen criterio de
evaluación; y cuando el conjunto no presenta contradicciones graves, duplicidades
evidentes ni vacíos estructurales.

\section{Validación de requisitos}
La validación contrasta con evidencia si los requisitos representan
adecuadamente las necesidades, restricciones, contexto de uso y criterios de
éxito del sistema; su pregunta central es si se está definiendo el sistema
correcto~\cite{iso29148}. Se diferencia de la verificación: esta comprueba si el
requisito está bien formulado, mientras que la validación comprueba si
corresponde a lo que realmente se necesita. En procesos orientados al
entendimiento compartido, la validación es crítica porque convierte la
interpretación documentada en un acuerdo explícito o en un desacuerdo gestionado.
La RSL advierte que la validación puede volverse superficial cuando se enfoca en la
forma textual y no en la cadena necesidad--requisito--criterio (RQ4, RQ5); la
caracterización adopta esa distinción como criterio de salida.

\subsection{Tareas, técnicas y criterios de salida}
La Tabla~\ref{tab:tareas_validacion} resume las tareas características. Como
técnicas se emplean revisiones con stakeholders clave, \emph{checklists} de
completitud, consistencia, viabilidad y testabilidad, prototipos de baja
fidelidad, \emph{mockups} navegables, diagramas de proceso, casos de uso y, en
requisitos críticos, simulación o modelado formal. La validación se considera
cerrada cuando los stakeholders relevantes revisaron el conjunto en un formato
comprensible, los defectos críticos están resueltos o registrados, existe
evidencia de acuerdo o aceptación condicionada, y la nueva versión queda trazable
y controlada.

\begin{table}[htbp]
\caption{Tareas características de la validación.}
\label{tab:tareas_validacion}
\centering
\footnotesize
\renewcommand{\arraystretch}{1.12}
\begin{tabularx}{\textwidth}{@{}L{4.0cm}Y L{3.2cm}@{}}
\toprule
\textbf{Tarea} & \textbf{Descripción operativa} & \textbf{Evidencia esperada} \\
\midrule
Preparar la validación & Definir participantes, alcance, criterios y tipo de
evidencia. & Plan o agenda de validación. \\
Retroalimentar requisitos & Devolver los requisitos para confirmar
interpretación. & Comentarios y observaciones. \\
Ejecutar revisión estructurada & \emph{Walkthroughs}, inspecciones, talleres o
\emph{checklists}. & Acta de revisión. \\
Usar artefactos de comprensión & Prototipos, modelos, escenarios o simulaciones.
& Evidencia revisada con stakeholders. \\
Evaluar contra contexto & Revisar si sirven para usuarios, tareas y ambientes
previstos. & Hallazgos por escenario. \\
Resolver conflictos & Negociar \emph{trade-offs} de alcance, costo, riesgo y
desempeño. & Decisiones y compromisos. \\
Obtener acuerdo & Registrar aprobación, rechazo o aceptación condicionada. &
Estado de aceptación por requisito. \\
Fijar baseline y registrar cambios & Consolidar versión controlada y documentar
ajustes. & Baseline versionada y registro de cambios. \\
\bottomrule
\end{tabularx}
\end{table}

\section{Elementos transversales}
Las tres actividades comparten elementos que sostienen el entendimiento
compartido al pasar de una fase a otra (Tabla~\ref{tab:transversales}). La
evidencia industrial muestra que una buena práctica de requisitos repercute en
productividad, calidad, planificación, gestión de riesgos y coordinación con
procesos posteriores~\cite{damian2006empirical}.

\begin{table}[htbp]
\caption{Elementos transversales a las tres actividades.}
\label{tab:transversales}
\centering
\footnotesize
\renewcommand{\arraystretch}{1.12}
\begin{tabularx}{\textwidth}{@{}L{4.0cm}Y@{}}
\toprule
\textbf{Elemento transversal} & \textbf{Función dentro del proceso} \\
\midrule
Gestión de stakeholders & Mantener visibles actores, responsabilidades, intereses
y capacidad de decisión. \\
Contexto de uso & Anclar requisitos en usuarios, tareas, ambientes y objetivos
reales. \\
Trazabilidad & Relacionar necesidades, requisitos, decisiones, cambios y
criterios de evaluación. \\
\emph{Rationale} & Explicar por qué se acepta, rechaza o modifica un requisito. \\
Calidad del lenguaje & Reducir ambigüedad y variación de interpretación. \\
Negociación & Gestionar conflictos entre stakeholders, restricciones y
prioridades. \\
Iteración & Permitir ajustes sin perder control del avance. \\
Versionado y baseline & Proteger los acuerdos logrados y controlar cambios. \\
Evidencia de entendimiento & Registrar acuerdos, aclaraciones, conflictos,
cobertura y retrabajo. \\
\bottomrule
\end{tabularx}
\end{table}

%=======================================================================
%  CAPÍTULO 4 — SME Y DECISIONES DE ADAPTACIÓN
%=======================================================================
\chapter{Ingeniería de métodos situacional y adaptación de THUNDERS}
\label{ch:adaptacion}

\section{Generalidades}
La adaptación de THUNDERS hacia THUNDERS-Master se fundamenta en la Ingeniería
de Métodos Situacionales (SME). Este capítulo describe SME, compara enfoques de
adaptación de procesos, justifica el método seleccionado y presenta las
decisiones de adaptación y su trazabilidad con THUNDERS.

\section{Ingeniería de Métodos Situacionales (SME)}
La SME propone construir métodos específicos para una situación a partir de
componentes reutilizables de métodos existentes~\cite{brinkkemper1996method,
henderson2010sme}. Incluye especificar los requisitos del método, seleccionar
componentes, ensamblarlos y validar el método resultante. Los \emph{method
fragments} y los \emph{method chunks} permiten tratar un proceso no como un
bloque indivisible, sino como una colección de piezas reutilizables; cada chunk
asocia una intención, una situación de aplicación, una guía de uso y productos
de trabajo~\cite{ralyte2001assembly,mirbel2006sme}. Este enfoque es
especialmente aplicable porque THUNDERS fue construido originalmente mediante
SME~\cite{agredo2022thesis}, lo que garantiza continuidad metodológica.

\section{Comparación de enfoques de adaptación}
Se compararon cinco enfoques —SME orientada por chunks, method fragments/chunks
como técnica aislada, \emph{software process tailoring}, \emph{software process
lines} y adaptación ágil basada en contexto— frente a criterios de adaptación de
procesos existentes, compatibilidad con IR, compatibilidad ágil, reducción de
complejidad, conservación de elementos esenciales y facilidad de aplicación en un
trabajo de grado. La Tabla~\ref{tab:comparacion} resume el análisis cualitativo.

\begin{table}[htbp]
\caption{Comparación de métodos de adaptación frente al caso THUNDERS-Master.}
\label{tab:comparacion}
\centering
\footnotesize
\renewcommand{\arraystretch}{1.15}
\begin{tabularx}{\textwidth}{@{}L{3.4cm}Y L{2.2cm}@{}}
\toprule
\textbf{Método} & \textbf{Valoración resumida} & \textbf{Referencia} \\
\midrule
SME orientada por chunks &
Alta capacidad de adaptar, orientar a elicitación/especificación/validación,
conservar intención y seleccionar solo componentes necesarios. Viable en trabajo
de grado. &
\cite{brinkkemper1996method} \\
\midrule
Method fragments/chunks (aislado) &
Alta reutilización y conservación de intención, pero requiere reglas de
composición y estructura de repositorio. &
\cite{ralyte2001assembly} \\
\midrule
Software process tailoring &
Buenos criterios de contexto, pero por sí solo no preserva la arquitectura
conceptual del método original. &
\cite{kalus2013tailoring} \\
\midrule
Software process lines &
Útil para familias de variantes a largo plazo, pero introduce más
infraestructura de la necesaria para una primera adaptación. &
\cite{henderson2010sme} \\
\midrule
Adaptación ágil basada en contexto &
Aporta ligereza, iteración y feedback rápido, pero no garantiza conservar el
núcleo conceptual de THUNDERS. &
\cite{campanelli2015agile} \\
\bottomrule
\end{tabularx}
\end{table}

\section{Método seleccionado}
El método seleccionado es una \emph{SME orientada por method chunks},
complementada con criterios de \emph{tailoring} contextual~\cite{kalus2013tailoring}
y reducción de carga cognitiva. La selección se justifica porque: THUNDERS fue
construido con SME~\cite{agredo2022thesis}; los chunks permiten una granularidad
adecuada para seleccionar lo que aporta valor directo a IR; el \emph{tailoring}
aporta criterios para ajustar el proceso al contexto; y la teoría de carga
cognitiva respalda reducir elementos que aumentan el esfuerzo mental sin aportar
valor proporcional. Frente a las alternativas de la Tabla~\ref{tab:comparacion}, la
SME orientada por chunks es la única que satisface a la vez los criterios decisivos
para este trabajo: continuidad metodológica con la construcción original de
THUNDERS, granularidad para conservar solo los componentes con aporte directo a la
IR, preservación de la intención del proceso y viabilidad de aplicación en un
trabajo de grado; el \emph{software process tailoring} y las \emph{software process
lines} no preservan por sí solos la arquitectura conceptual del método, y la
adaptación ágil basada en contexto no garantiza conservar su núcleo. La
Tabla~\ref{tab:sme_application} muestra cómo se concretó SME en el diseño.

\begin{table}[htbp]
\caption{Aplicación de las actividades de SME en el diseño de THUNDERS-Master.}
\label{tab:sme_application}
\centering
\footnotesize
\renewcommand{\arraystretch}{1.15}
\begin{tabularx}{\textwidth}{@{}L{2.9cm}Y Y@{}}
\toprule
\textbf{Actividad SME} & \textbf{Aplicación en THUNDERS-Master} &
\textbf{Salida de diseño} \\
\midrule
Especificar requisitos del método &
Definir qué debe permitir el proceso: construir entendimiento compartido,
reducir ambigüedad, registrar decisiones, validar criterios y mantener
trazabilidad mínima. &
Matriz de requisitos del método, conjunto de evidencia mínima y criterios de
verificación. \\
\midrule
Seleccionar componentes &
Analizar actividades, roles y artefactos de THUNDERS según su aporte al
entendimiento compartido, relevancia para IR y costo de adopción. &
Matriz de decisión: conservar, adaptar, fusionar o eliminar. \\
\midrule
Ensamblar componentes &
Reorganizar los componentes en un flujo de tres fases compatible con sesiones
de requisitos. &
Modelo conceptual operativo con actividades A1--A11 y la opcional A4b. \\
\midrule
Revisar coherencia &
Separar la operación del proceso de su evaluación como investigación y verificar
trazabilidad entre requisitos, actividades y artefactos. &
Rúbrica de revisión experta, walkthrough, plan GQM y criterios de refinamiento. \\
\bottomrule
\end{tabularx}
\end{table}

\section{Decisiones de adaptación y trazabilidad con THUNDERS}
La regla principal es conservar la intención de THUNDERS, no su complejidad
completa: un componente solo pasa al núcleo operativo si aporta de forma directa
al entendimiento compartido en IR y si su salida puede comprobarse mediante
evidencia mínima. Así, la caracterización amplia de participantes se convierte en
un mapa mínimo de stakeholders; el diseño completo de actividad se simplifica como
canvas de sesión; la evaluación de desempeño de participantes se elimina del
núcleo operativo; y la validación de solución se reinterpreta como validación
semántica de la cadena necesidad--requisito--criterio.

\subsection{Criterios de decisión de adaptación}
Para que estas decisiones no dependan del juicio puntual de los autores, cada
componente de THUNDERS se evaluó con un conjunto explícito y estable de criterios,
anclados en la evidencia de la RSL (Sección~\ref{sec:rsl}) y en la caracterización
de la IR (Capítulo~\ref{ch:caracterizacion}). La Tabla~\ref{tab:decision_criteria}
define esos criterios y su fundamento; la decisión sobre cada componente
—conservar, adaptar, simplificar, fusionar, hacer opcional o eliminar— resulta de
aplicarlos de forma conjunta.

\begin{table}[htbp]
\caption{Criterios de decisión para adaptar los componentes de THUNDERS.}
\label{tab:decision_criteria}
\centering
\footnotesize
\renewcommand{\arraystretch}{1.15}
\begin{tabularx}{\textwidth}{@{}L{4.2cm}Y L{3.0cm}@{}}
\toprule
\textbf{Criterio} & \textbf{Pregunta que responde} & \textbf{Fundamento} \\
\midrule
Aporte directo al entendimiento compartido en IR & ¿El componente contribuye a
construir, monitorear o validar el significado compartido de los requisitos? &
Hallazgos de la RSL (RQ1--RQ4). \\
Verificabilidad por evidencia mínima & ¿Su salida puede comprobarse con un artefacto
o evidencia observable y ligera? & Indicadores indirectos de la RSL (RQ5). \\
Costo de adopción y carga cognitiva & ¿El esfuerzo que introduce es proporcional a
su valor en una sesión breve? & Validación de THUNDERS como proceso extenso y
demandante~\cite{agredo2022exploratory}. \\
Pertinencia para elicitación, especificación y validación & ¿El componente
corresponde al alcance de las tres actividades caracterizadas? & Caracterización de
la IR (Capítulo~\ref{ch:caracterizacion}). \\
\bottomrule
\end{tabularx}
\end{table}

La Tabla~\ref{tab:thunders_traceability} aplica estos criterios y resume la decisión
adoptada para cada componente.

\begin{table}[htbp]
\caption{Trazabilidad resumida entre THUNDERS y THUNDERS-Master.}
\label{tab:thunders_traceability}
\centering
\footnotesize
\renewcommand{\arraystretch}{1.12}
\begin{tabularx}{\textwidth}{@{}L{3.6cm}L{2.6cm}Y@{}}
\toprule
\textbf{Componente THUNDERS} & \textbf{Decisión} & \textbf{Justificación} \\
\midrule
C-01 Definir la población & Adaptar &
La caracterización completa es costosa; en IR basta con asegurar
representación, conocimiento y autoridad de decisión. \\
C-02 Definir el tema de la actividad & Conservar/Adaptar &
Representa el punto inicial de alineación entre problema, alcance y objetivo. \\
C-03 Diseñar la actividad & Simplificar/Fusionar &
Un diseño completo es demasiado pesado para sesiones ágiles de requisitos. \\
C-04 Definir los grupos & Opcional (activable) &
No siempre se requieren múltiples grupos; se prioriza participación
representativa. \\
C-05 Diseño del material & Simplificar &
Preparar material es útil, pero no debe duplicar documentación existente. \\
C-06 Diseñar validación y evaluación & Dividir &
Separa la validación semántica de requisitos de la evaluación del proceso como
investigación. \\
C-07 Formar los grupos & Eliminar/Fusionar &
Duplica la definición de participación y aporta poco en sesiones pequeñas. \\
C-08 Describir la actividad & Conservar/Adaptar &
Momento clave para confirmar el entendimiento inicial del problema. \\
C-09 Distribuir el material & Fusionar &
No debe ser actividad independiente cuando los artefactos ya están disponibles. \\
C-10 Desarrollar la actividad colaborativa & Adaptar como núcleo &
Se convierte en la actividad central, transformada al flujo real de IR. \\
C-11 Validar la solución & Adaptar como núcleo &
Se enfoca en validar la correspondencia semántica, no una solución
implementada. \\
C-12 Evaluar desempeño de participantes & Eliminar del núcleo &
No mide directamente calidad de requisitos ni entendimiento compartido y
aumenta la sobrecarga. \\
C-13 Cerrar la actividad & Conservar/Adaptar &
Evita la pérdida de contexto y asegura que los acuerdos sigan siendo
trazables. \\
\bottomrule
\end{tabularx}
\end{table}

El resultado de aplicar estos criterios es un núcleo deliberadamente ligero, cuya
baja carga es la condición que permite integrar THUNDERS-Master dentro de flujos
ágiles sin introducir burocracia (Sección~\ref{sec:encaje_agil},
Tabla~\ref{tab:agile_fit}).

%=======================================================================
%  CAPÍTULO 5 — THUNDERS-MASTER COMO PROCESO OPERATIVO (LAS ETAPAS)
%=======================================================================
\chapter{THUNDERS-Master como proceso operativo}
\label{ch:proceso}

\notaversion{La especificación presentada en este capítulo corresponde a la
\textbf{primera versión de THUNDERS-Master (v1.0)}: una especificación formal de
trabajo derivada de la RSL, la caracterización de la IR, el análisis de THUNDERS
y la Ingeniería de Métodos Situacionales. No es un proceso empíricamente
validado. Está concebida para evolucionar de manera iterativa e incremental
mediante un taller de validación con expertos y usuarios clave, y un pilotaje en
un contexto real (agrícola). La retroalimentación obtenida producirá versiones
posteriores (v1.1, v2.0, \dots) cuyos cambios quedarán documentados. Por ello,
fases, actividades, roles y artefactos deben leerse como una base estable pero
abierta a refinamiento, no como un diseño definitivo.}

\section{Generalidades}
THUNDERS-Master es un proceso ligero para apoyar la construcción, el monitoreo y
la validación del entendimiento compartido durante la elicitación,
especificación y validación de requisitos. No reemplaza Scrum, Kanban, Design
Thinking, refinamiento de backlog ni otras prácticas existentes: las complementa
con orientación explícita para hacer visible el significado común y dejar
evidencia mínima que permita seguir trabajando sin perder contexto. La unidad de
aplicación es una sesión, taller, refinamiento o revisión de requisitos. El
proceso se estructura en tres fases heredadas de THUNDERS y reinterpretadas para
IR —Beginning, Developing y Measuring/Closing— y se desarrolla mediante once
actividades operativas (A1--A11) más una actividad opcional (A4b). La
Figura~\ref{fig:modelo} presenta el modelo conceptual operativo.

\begin{figure}[htbp]
\centering
\begin{tikzpicture}[
  font=\footnotesize,
  fase/.style={draw,rounded corners=3pt,align=center,minimum width=3.6cm,
               minimum height=1.0cm,fill=blue!10,thick},
  act/.style={draw,rounded corners=2pt,align=left,text width=3.5cm,
              minimum height=0.8cm,fill=gray!5,font=\scriptsize},
  arr/.style={-{Stealth[length=2.4mm]},thick}]

  % fases
  \node[fase] (f1) {\textbf{Fase 1}\\Beginning / Preparar};
  \node[fase,right=1.2cm of f1] (f2) {\textbf{Fase 2}\\Developing / Construir};
  \node[fase,right=1.2cm of f2] (f3) {\textbf{Fase 3}\\Measuring / Validar y cerrar};
  \draw[arr] (f1)--(f2);
  \draw[arr] (f2)--(f3);

  % actividades bajo cada fase
  \node[act,below=0.5cm of f1] (a1) {A1 Canvas de sesión\\A2 Stakeholders mínimos\\A3 Problema + glosario\\A4 Materiales mínimos\\\textit{A4b Subgrupos (opc.)}};
  \node[act,below=0.5cm of f2] (a2) {A5 Apertura de alineación\\A6 Necesidades y supuestos\\A7 Requisitos + criterios\\A8 Chequeo de entend. + traz.};
  \node[act,below=0.5cm of f3] (a3) {A9 Validar requisitos\\A10 Resolver discrepancias\\A11 Baseline y acciones};
  \draw[arr] (f1)--(a1);
  \draw[arr] (f2)--(a2);
  \draw[arr] (f3)--(a3);

  % línea de evidencia mínima
  \node[draw,fill=green!8,rounded corners=2pt,below=0.5cm of a2,
        text width=11.8cm,align=center,font=\scriptsize] (ev)
        {Línea de evidencia mínima: canvas, mapa de stakeholders, glosario,
         registro de necesidades, tarjetas de requisito, matriz de trazabilidad,
         decision log y baseline.};
\end{tikzpicture}
\caption{Modelo conceptual operativo de THUNDERS-Master.}
\label{fig:modelo}
\end{figure}

El modelo expresa una idea clave: el entendimiento compartido suficiente no se
asume por la existencia de documentos. Se considera construido cuando los
participantes pueden reformular el problema de forma equivalente, los términos
críticos tienen definición acordada o responsable, cada requisito crítico
mantiene relación con una necesidad y un criterio, los desacuerdos quedan
resueltos o asignados, y el cierre produce una versión identificable de los
requisitos trabajados.

\section{Encaje de THUNDERS-Master en marcos ágiles}
\label{sec:encaje_agil}
El propósito de THUNDERS-Master es apoyar la construcción, el monitoreo y la
validación del entendimiento compartido; por eso se concibe como una buena práctica
que se aplica \emph{dentro} de las metodologías ágiles, no como un marco que las
reemplace. Marcos como Scrum o Kanban ya coordinan el trabajo mediante ceremonias,
refinamiento y revisiones, pero no prescriben cómo asegurar que los participantes
interpreten de la misma manera el problema, las necesidades y los criterios de
aceptación. THUNDERS-Master cubre precisamente esa capa semántica y puede integrarse
en los espacios ágiles existentes sin crear reuniones nuevas: sus actividades se
acoplan a ceremonias que el equipo ya realiza (Tabla~\ref{tab:agile_fit}).

Este encaje también responde a la preocupación por una posible sobrecarga
documental. Los artefactos de THUNDERS-Master no son entregables adicionales, sino
campos mínimos que viven dentro de las herramientas que el equipo ya usa —backlog,
wiki, \emph{issues}, actas o documentos colaborativos (Sección~\ref{sec:artefactos})—
y que solo se intensifican cuando el contexto lo exige (Tabla~\ref{tab:activacion}).
Así, el proceso agrega trazabilidad y evidencia de entendimiento sin penalizar la
agilidad: la documentación se limita a lo que permite verificar quién originó un
requisito, qué necesidad expresa y qué criterio demuestra que se cumple.

\begin{table}[htbp]
\caption{Encaje de las actividades de THUNDERS-Master en ceremonias ágiles.}
\label{tab:agile_fit}
\centering
\footnotesize
\renewcommand{\arraystretch}{1.15}
\begin{tabularx}{\textwidth}{@{}L{3.2cm}L{3.8cm}Y@{}}
\toprule
\textbf{Actividades T-M} & \textbf{Espacio ágil de acople} & \textbf{Aporte al
entendimiento compartido} \\
\midrule
A1--A4 (Preparar) & Preparación de \emph{sprint planning} y refinamiento;
\emph{kickoff} de iniciativa. & Fija objetivo, stakeholders, problema y glosario
antes de comprometer trabajo. \\
A5--A8 (Construir) & Refinamiento de \emph{backlog} y escritura de
historias/criterios. & Convierte necesidades en requisitos trazables y aplica
micro-chequeos de interpretación. \\
A9--A11 (Validar y cerrar) & \emph{Sprint review}, aceptación y \emph{Definition of
Done}. & Valida la cadena necesidad--requisito--criterio y fija una baseline
acordada. \\
\bottomrule
\end{tabularx}
\end{table}

\section{Requisitos del método}
Los requisitos del método funcionan como contrato de diseño: permiten comprobar
que el proceso no se limita a renombrar fases de THUNDERS, sino que responde a
necesidades específicas de la IR. Se consolidan requisitos relacionados con
stakeholders, contexto de sesión, problema compartido, lenguaje común,
necesidades y supuestos, transformación a requisitos, calidad semántica,
trazabilidad mínima, criterios de aceptación, decisiones y desacuerdos, chequeo
de entendimiento, asistencia ligera, variabilidad contextual, evaluación externa
y no transferencia excesiva. El proceso distingue elementos obligatorios
(protegen el significado mínimo), recomendados (ayudan cuando la sesión es
compleja) y opcionales (se activan por contexto: distribución, criticidad,
regulación o baja madurez del equipo).

\section{Fase 1: Beginning / Preparar}
La fase \emph{Beginning} prepara la sesión antes de elicitar o especificar. Su
propósito es reducir incertidumbre inicial, delimitar el trabajo y asegurar que
participen las fuentes correctas.
\begin{itemize}
  \item \act{A1 Canvas de sesión.} Define objetivo, alcance, exclusiones,
        modalidad, agenda, técnica, insumos, salida esperada, regla de
        participación y criterio de cierre.
  \item \act{A2 Stakeholders mínimos.} Identifica actores, fuente de
        conocimiento, poder de decisión, disponibilidad, posibles conflictos y
        prioridad de participación.
  \item \act{A3 Problema y glosario semilla.} Formula una declaración compartida
        del problema y registra términos de dominio que pueden generar
        ambigüedad.
  \item \act{A4 Materiales y validación mínima.} Prepara insumos, preguntas guía,
        plantillas y criterios que se usarán durante la sesión.
  \item \act{A4b Subgrupos/turnos (opcional).} Se activa si hay muchos
        participantes, asimetrías de poder, distribución geográfica o necesidad
        de asegurar participación equilibrada.
\end{itemize}

\section{Fase 2: Developing / Construir}
La fase \emph{Developing} es el núcleo del proceso. Su objetivo es construir
requisitos de manera colaborativa mientras se monitorea la comprensión
compartida.
\begin{itemize}
  \item \act{A5 Apertura de alineación.} Revisa objetivo, alcance, reglas,
        participantes, problema compartido y términos críticos para partir de una
        base semántica común.
  \item \act{A6 Necesidades y supuestos.} Elicita necesidades, restricciones,
        supuestos, valor esperado y contexto de uso; cada necesidad debe tener
        fuente, contexto y estado.
  \item \act{A7 Requisitos/historias y criterios.} Transforma necesidades en
        requisitos, historias, escenarios o criterios de aceptación, conservando
        la relación con la necesidad fuente.
  \item \act{A8 Chequeo de entendimiento y trazabilidad.} Aplica
        \emph{micro-chequeos} en puntos críticos: pedir reformulación, contrastar
        interpretación, revisar criterios observables y completar una matriz
        mínima de trazabilidad.
\end{itemize}
Los micro-chequeos son una decisión central: THUNDERS-Master no espera al final
para descubrir discrepancias, sino que introduce revisiones breves durante la
construcción (por ejemplo, pedir a dos participantes que reformulen un requisito,
preguntar qué evidencia demostraría que un criterio se cumple, o marcar un
supuesto como pendiente cuando falta información).

\section{Fase 3: Measuring / Validar y cerrar}
La fase \emph{Measuring} se reinterpreta para IR: la evaluación de desempeño de
participantes no pertenece al núcleo operativo. La fase se concentra en validar
el significado de los requisitos, resolver discrepancias y cerrar una baseline.
\begin{itemize}
  \item \act{A9 Validar requisitos y criterios.} Revisa la cadena
        necesidad--requisito--criterio--significado con los stakeholders. Un
        requisito no se considera listo solo por estar bien escrito: debe estar
        vinculado a una necesidad, contar con criterio observable y ser
        interpretado coherentemente.
  \item \act{A10 Resolver discrepancias.} Registra diferencias de interpretación,
        conflictos de alcance, desacuerdos de prioridad o criterios no
        observables; cada discrepancia se resuelve o queda asignada con
        responsable y fecha.
  \item \act{A11 Baseline y acciones.} Cierra la sesión con una versión
        identificable de los requisitos, decision log, pendientes, responsables y
        próximos pasos.
\end{itemize}

La Tabla~\ref{tab:actividades} sintetiza las once actividades con su entrada,
salida y criterio de cierre.

\begin{table}[htbp]
\caption{Especificación resumida de las actividades de THUNDERS-Master.}
\label{tab:actividades}
\centering
\footnotesize
\renewcommand{\arraystretch}{1.12}
\begin{tabularx}{\textwidth}{@{}L{1.5cm}L{2.7cm}Y L{2.8cm}@{}}
\toprule
\textbf{Actividad} & \textbf{Propósito} & \textbf{Salida / evidencia mínima} &
\textbf{Criterio de cierre} \\
\midrule
A1 Canvas & Definir objetivo, alcance, agenda y cierre. & Canvas breve de sesión.
& La sesión tiene objetivo y salida esperada. \\
A2 Stakeholders & Asegurar fuentes de conocimiento y autoridad. & Mapa mínimo de
stakeholders. & Los requisitos críticos tienen fuente o validador. \\
A3 Problema + glosario & Alinear el problema y la terminología crítica. &
Enunciado del problema y glosario semilla. & Los participantes reformulan el
problema de forma consistente. \\
A4 Materiales & Preparar insumos y plantillas. & Paquete mínimo de artefactos
referenciados. & Los insumos están disponibles sin duplicar documentación. \\
A5 Apertura & Confirmar un punto de partida común. & Ajustes iniciales y reglas
acordadas. & No quedan dudas críticas sobre el alcance. \\
A6 Necesidades & Elicitar necesidades con fuente y contexto. & Registro de
necesidades y supuestos. & Cada necesidad tiene fuente, contexto y estado. \\
A7 Requisitos + criterios & Transformar necesidades en artefactos. & Requisitos o
historias y criterios de aceptación. & Cada requisito se enlaza con su necesidad
de origen. \\
A8 Chequeo + trazabilidad & Monitorear interpretación y trazabilidad. & Checklist
y matriz mínima de trazabilidad. & Las discrepancias críticas están
identificadas. \\
A9 Validar & Confirmar correspondencia y acuerdo. & Requisitos aceptados o
corregidos. & Los criterios son observables y validados. \\
A10 Discrepancias & Gestionar diferencias de significado. & Decision log y
acciones. & Todo desacuerdo crítico se resuelve o asigna. \\
A11 Baseline & Establecer baseline y continuidad. & Baseline, responsables y
próximos pasos. & Existe una versión identificable de lo trabajado. \\
\bottomrule
\end{tabularx}
\end{table}

\section{Roles y responsabilidades}
Los roles se reducen para evitar especialización innecesaria. En equipos
pequeños, una persona puede asumir más de un rol, siempre que no se pierdan tres
responsabilidades: facilitar la interacción, especificar requisitos y registrar
evidencias. La Tabla~\ref{tab:roles} resume los roles simplificados.

\begin{table}[htbp]
\caption{Roles simplificados en THUNDERS-Master.}
\label{tab:roles}
\centering
\footnotesize
\renewcommand{\arraystretch}{1.12}
\begin{tabularx}{\textwidth}{@{}L{2.9cm}Y L{2.6cm}@{}}
\toprule
\textbf{Rol} & \textbf{Responsabilidad principal} & \textbf{Condición de uso} \\
\midrule
Facilitador de requisitos & Conducir la sesión, gestionar participación, aplicar
micro-chequeos y cerrar. & Obligatorio. \\
Analista / especificador & Transformar necesidades en requisitos, historias,
escenarios y criterios. & Obligatorio. \\
Relator / gestor de trazabilidad & Registrar términos, necesidades, decisiones,
desacuerdos, enlaces y baseline. & Recomendado; obligatorio en sesiones grandes. \\
Stakeholder de dominio & Expresar necesidades, validar significados y aclarar
terminología. & Obligatorio. \\
Responsable de decisión / PO & Priorizar, confirmar alcance y aprobar criterios.
& Obligatorio cuando el rol existe en el contexto. \\
Representante técnico / QA & Evaluar factibilidad, verificabilidad y testabilidad.
& Recomendado. \\
Equipo investigador externo & Evaluar utilidad, claridad, carga y calidad de
artefactos. & Externo al proceso operativo. \\
\bottomrule
\end{tabularx}
\end{table}

\section{Artefactos mínimos y evidencia observable}
\label{sec:artefactos}
Los artefactos se diseñan como campos mínimos que pueden vivir dentro de
herramientas existentes (backlog, wiki, issues, actas o documentos
colaborativos): canvas de sesión, mapa de stakeholders, glosario vivo, registro
de necesidades, tarjeta de requisito o historia, matriz mínima de trazabilidad,
decision log y baseline. La evidencia mínima no busca llenar formatos por
obligación, sino permitir verificación posterior: un revisor experto debería
poder responder quién originó un requisito, qué necesidad expresa, qué criterio
demuestra que se cumple, qué término crítico se acordó, qué desacuerdo quedó
abierto, quién decide y qué versión fue validada.

\section{Variabilidad contextual y criterios de activación}
THUNDERS-Master no se ejecuta siempre con la misma intensidad. La variabilidad
contextual define cuándo activar elementos adicionales, en línea con la
adaptación estructurada de procesos~\cite{kalus2013tailoring} y evitando su
aplicación rígida en contextos ágiles~\cite{campanelli2015agile}. La
Tabla~\ref{tab:activacion} resume los criterios.

\begin{table}[htbp]
\caption{Criterios de activación contextual en THUNDERS-Master.}
\label{tab:activacion}
\centering
\footnotesize
\renewcommand{\arraystretch}{1.12}
\begin{tabularx}{\textwidth}{@{}L{4.2cm}Y Y@{}}
\toprule
\textbf{Criterio de activación} & \textbf{Elemento activado} &
\textbf{Evidencia esperada} \\
\midrule
Alta heterogeneidad de stakeholders & Subgrupos, rondas de participación o
reformulación cruzada. & Registro de perspectivas y acuerdos integrados. \\
Requisitos críticos o regulados & Trazabilidad extendida y validación más formal.
& Enlaces completos necesidad--requisito--criterio--decisión--evidencia. \\
Distribución geográfica o comunicación asíncrona & Canvas reforzado, glosario vivo
y decision log más explícito. & Acuerdos visibles para quienes no asistieron. \\
Baja madurez en IR & Plantillas de una página, checklist de calidad y preguntas
guía. & Requisitos con estructura mínima y criterios observables. \\
Terminología de dominio ambigua & Glosario vivo y responsable de resolución
semántica. & Términos críticos definidos o marcados como pendientes. \\
Discrepancias no resueltas & Sesión adicional o proceso de escalamiento. &
Discrepancia asignada con responsable, fecha y criterio de cierre. \\
\bottomrule
\end{tabularx}
\end{table}

\section{Estado de madurez y evolución del proceso}
THUNDERS-Master se asume como una especificación de trabajo en evolución. Su
construcción y refinamiento siguen la lógica iterativa e incremental con que se
construyó THUNDERS, sometido a sucesivas validaciones que corregían y
enriquecían cada versión~\cite{agredo2022thesis,agredo2022exploratory}. La
Tabla~\ref{tab:versionado} describe la ruta de evolución prevista: la versión
v1.0 (la especificada en este capítulo) alimenta un ciclo de validación experta y
pilotaje en contexto agrícola, del cual se derivarán versiones posteriores. Cada
iteración debe registrar qué cambió, por qué cambió y con qué evidencia, de modo
que la trazabilidad entre versiones se conserve.

\begin{table}[htbp]
\caption{Ruta de evolución prevista para THUNDERS-Master.}
\label{tab:versionado}
\centering
\footnotesize
\renewcommand{\arraystretch}{1.12}
\begin{tabularx}{\textwidth}{@{}L{2.2cm}L{4.6cm}Y@{}}
\toprule
\textbf{Versión} & \textbf{Entrada / fuente de cambios} & \textbf{Resultado
esperado} \\
\midrule
v1.0 & RSL, caracterización, análisis de THUNDERS y SME. & Especificación formal
de trabajo (este capítulo). \\
v1.1 & Taller de validación con expertos y usuarios clave (coherencia interna,
pertinencia de actividades, claridad operativa). & Versión refinada lista para
pilotar, con ajustes documentados. \\
v2.0 & Pilotaje en estudio de caso agrícola y ejecución del plan GQM (utilidad,
facilidad de uso, carga cognitiva, calidad de artefactos). & Versión ajustada con
lecciones aprendidas y evidencia empírica. \\
Futuras & Réplicas en otros dominios y contextos. & Posible línea de procesos
(SPrL) con variantes especializadas. \\
\bottomrule
\end{tabularx}
\end{table}

\notaversion{Los resultados de las iteraciones v1.1 y v2.0 se reportarán en el
Capítulo~\ref{ch:evaluacion}. Mientras esas validaciones no se ejecuten, las
afirmaciones sobre utilidad y facilidad de uso de THUNDERS-Master deben
entenderse como hipótesis de diseño, no como hallazgos confirmados.}

%=======================================================================
%  CAPÍTULO 6 — EVALUACIÓN (ETAPAS VACÍAS / TRABAJO FUTURO)
%=======================================================================
\chapter{Evaluación de THUNDERS-Master}
\label{ch:evaluacion}

\section{Generalidades}
La evaluación se separa del uso operativo del proceso: durante una sesión real no
se incorpora un rol investigador que evalúe desempeño individual ni se aplican
instrumentos extensos. La evaluación pertenece al proyecto de investigación y se
orienta a validar la coherencia, utilidad, facilidad de uso y capacidad del
proceso para producir evidencia de entendimiento compartido.

\section{Plan de evaluación basado en GQM}
El plan inicial se estructura mediante Goal-Question-Metric
(GQM)~\cite{basili1994gqm}. El objetivo es evaluar si THUNDERS-Master ayuda a
construir, monitorear y validar entendimiento compartido durante la elicitación,
especificación y validación. A partir de este objetivo se formulan preguntas
sobre reducción de ambigüedad, trazabilidad, discrepancias, calidad de criterios,
utilidad percibida, facilidad de aplicación y carga
cognitiva~\cite{hart1988nasa}. Las métricas se obtendrán mediante revisión de
artefactos, checklist de calidad, decision log, encuestas breves —la escala
SUS~\cite{brooke1996sus} puede servir como referencia de usabilidad—, entrevistas
y observación del piloto. La Tabla~\ref{tab:gqm} presenta la propuesta inicial.

\begin{table}[htbp]
\caption{Plan de evaluación propuesto para THUNDERS-Master (GQM).}
\label{tab:gqm}
\centering
\footnotesize
\renewcommand{\arraystretch}{1.12}
\begin{tabularx}{\textwidth}{@{}L{3.0cm}Y L{3.2cm}@{}}
\toprule
\textbf{Objetivo} & \textbf{Pregunta / indicador} & \textbf{Instrumento} \\
\midrule
Evaluar utilidad & ¿El proceso ayuda a construir entendimiento compartido?
Malentendidos detectados, acuerdos logrados. & Cuestionario y revisión de
artefactos. \\
Evaluar facilidad de uso & ¿Se aplica sin esfuerzo excesivo? Claridad de pasos,
tiempo de aplicación. & Cuestionario de usabilidad (SUS). \\
Evaluar carga cognitiva & ¿Reduce el esfuerzo mental frente a THUNDERS completo?
& NASA-TLX. \\
Evaluar calidad de artefactos & ¿Los requisitos son comprensibles, trazables y
validables? & Checklist de calidad de requisitos. \\
Evaluar conservación del propósito & ¿Mantiene construcción, monitoreo y
validación del entendimiento? & Matriz de trazabilidad entre THUNDERS y
THUNDERS-Master. \\
\bottomrule
\end{tabularx}
\end{table}

\section{Estrategia de evaluación}
La evaluación se ejecuta en dos instancias complementarias, alineadas con las
fases de la investigación (Sección~\ref{sec:metodologia}) y con la ruta de
versionado del proceso (Tabla~\ref{tab:versionado}). La primera es un
\emph{taller de validación con expertos y usuarios clave}, orientado a verificar
la coherencia interna de THUNDERS-Master, la pertinencia de sus actividades,
roles e instrumentos, y la claridad operativa de sus pasos; su retroalimentación
produce la versión a pilotar (v1.1) y permite revisar la alineación entre el
proceso y el plan de medición GQM. La segunda es un \emph{estudio de caso} que
aplica el proceso en un contexto real (agrícola) y ejecuta el plan GQM para
obtener evidencia empírica (v2.0). Esta separación evita confundir la ejecución
operativa del proceso con su evaluación como investigación.

\section{Diseño del estudio de caso en contexto agrícola}
El pilotaje se realizará en un contexto agrícola por las razones expuestas en la
Sección~\ref{sec:contexto_agricola}: diversidad técnica, disciplinar y
comunicativa entre productores, asesores, técnicos, investigadores y diseñadores,
que incrementa el riesgo de malentendidos y hace especialmente pertinente un
proceso de alineación semántica~\cite{getson2022communication,sulaiman2021ict}.
Tras seleccionar y formalizar el caso, la implementación se organiza en sesiones
que corresponden a las tres actividades del alcance: \emph{(i)}~elicitación, con
sesiones de levantamiento para capturar necesidades, restricciones y criterios de
aceptación; \emph{(ii)}~especificación, estructurando lo capturado en artefactos
acordados (por ejemplo, historias de usuario o especificación de requisitos)
mediante las guías y plantillas de THUNDERS-Master; y \emph{(iii)}~validación, con
sesiones de revisión y acuerdo que culminan en una versión acordada del conjunto
de requisitos.

Durante la evaluación se ejecuta el plan GQM con los instrumentos definidos. El
análisis combina estadística descriptiva para las métricas de percepción y
análisis temático para la información cualitativa de entrevistas y observaciones,
con el fin de identificar patrones de entendimiento compartido y puntos de
fricción. Con base en este análisis se elabora un registro de lecciones
aprendidas y se definen los ajustes a THUNDERS-Master.

\section{Validación por expertos}
\pendiente{Reportar la ejecución del taller de validación con expertos y usuarios
clave: perfil y número de participantes, instrumento (rúbrica de completitud,
coherencia, claridad, aplicabilidad y alineación con IR), protocolo de aplicación,
hallazgos y ajustes realizados a la especificación (v1.0~$\rightarrow$~v1.1).}

\section{Walkthrough metodológico}
\pendiente{Realizar un recorrido (walkthrough) paso a paso del proceso sobre un
caso representativo para verificar coherencia interna, dependencias entre
actividades y completitud de artefactos. Documentar inconsistencias y
correcciones.}

\section{Resultados del pilotaje en contexto agrícola}
\pendiente{Presentar el desarrollo del estudio de caso: contexto y participantes,
configuración del proceso (núcleo y elementos activados), ejecución de las
sesiones de elicitación, especificación y validación, y recolección de evidencia
según el plan GQM. Incluir las amenazas a la validez del estudio.}

\section{Análisis de resultados y lecciones aprendidas}
\pendiente{Presentar y discutir los resultados (utilidad, facilidad de uso, carga
cognitiva, calidad de artefactos y conservación del propósito), compararlos con
las expectativas del plan GQM, consolidar el registro de lecciones aprendidas y
derivar la versión ajustada de THUNDERS-Master (v2.0).}

%=======================================================================
%  CAPÍTULO 7 — CONCLUSIONES Y TRABAJO FUTURO
%=======================================================================
\chapter{Conclusiones y trabajo futuro}
\label{ch:conclusiones}

\section{Generalidades}
Este capítulo recoge las conclusiones del trabajo, sus contribuciones,
limitaciones y las líneas de trabajo futuro.

\section{Conclusiones}
El entendimiento compartido es una condición crítica para la calidad de los
requisitos, porque permite que stakeholders y equipos técnicos interpreten de
forma coherente problemas, necesidades, restricciones y criterios de validación.
La RSL muestra que las rupturas de significado aparecen en la identificación de
stakeholders, el lenguaje del dominio, la captura de necesidades, la
especificación, la trazabilidad y la validación. Sobre esa base, la revisión
delimitó una brecha concreta (Sección~\ref{sec:brecha}): no existe una ruta
integrada, ligera y trazable para construir, monitorear y validar el entendimiento
compartido, lo que justifica adaptar THUNDERS en lugar de crear un proceso nuevo. A
partir de esa evidencia y de la caracterización de las tres actividades, este
trabajo definió THUNDERS-Master:
un proceso ligero, derivado de THUNDERS mediante SME, para construir, monitorear
y validar el entendimiento compartido en IR. La especificación formaliza
requisitos del método, decisiones de adaptación, tres fases, once actividades
operativas, roles simplificados, artefactos mínimos, criterios de activación
contextual y un plan inicial de evaluación.
\pendiente{Completar las conclusiones con los resultados de la evaluación
(Capítulo~\ref{ch:evaluacion}) una vez ejecutada la validación experta y el
piloto.}

\section{Contribuciones}
Las contribuciones del trabajo, alineadas con la Sección~\ref{sec:metodologia} y
los aportes descritos en el Capítulo~\ref{ch:introduccion}, son de tres tipos.
En el ámbito \emph{académico}, amplía la comprensión teórica y práctica del
entendimiento compartido en IR, ofreciendo una definición operacional y una
trazabilidad explícita entre THUNDERS y THUNDERS-Master. En el ámbito
\emph{práctico/industrial}, aporta un proceso formal pero liviano —con fases,
actividades A1--A11, roles simplificados, artefactos mínimos y criterios de
activación contextual— que busca reducir ambigüedades, retrabajo y errores de
interpretación. En el ámbito \emph{investigativo}, extiende la aplicación de
THUNDERS al dominio de la IR y prepara la generación de evidencia empírica en un
contexto agrícola.
\pendiente{Consolidar las contribuciones con la evidencia obtenida tras la
validación experta y el pilotaje (Capítulo~\ref{ch:evaluacion}).}

\section{Limitaciones}
La principal limitación es que esta versión sigue siendo una especificación
formal de trabajo. Aunque está fundamentada en la RSL, la caracterización de IR,
el análisis de THUNDERS y SME, todavía no debe presentarse como validada
empíricamente.
\pendiente{Ampliar las limitaciones con las amenazas a la validez identificadas
durante la evaluación.}

\section{Trabajo futuro}
El trabajo futuro inmediato consiste en ejecutar la revisión experta, el
walkthrough metodológico y un piloto o estudio de caso, analizando claridad,
utilidad, facilidad de aplicación, carga percibida, calidad de artefactos y
señales de entendimiento compartido. A mediano plazo se plantea construir un
repositorio de \emph{method chunks} THUNDERS-RE y explorar una línea de procesos
(SPrL) que permita derivar variantes para otros dominios (educación, salud,
contextos empresariales).

\section{Productos de investigación}
\pendiente{Listar los productos derivados del trabajo: artículos enviados o
publicados, ponencias y demás actividades de investigación.}

%=======================================================================
%  BIBLIOGRAFÍA
%=======================================================================
\begin{thebibliography}{99}
\addcontentsline{toc}{chapter}{Bibliografía}

\bibitem{iso29148} ISO/IEC/IEEE, \emph{ISO/IEC/IEEE 29148:2018 — Systems and
software engineering — Life cycle processes — Requirements engineering}, 2nd ed.,
2018.

\bibitem{iso25010} ISO/IEC, \emph{ISO/IEC 25010:2023 — Systems and software
engineering — SQuaRE — Product quality model}, 2nd ed., 2023.

\bibitem{iso9241} ISO, \emph{ISO 9241-11:2018 — Ergonomics of human-system
interaction — Part 11: Usability: Definitions and concepts}, 2nd ed., 2018.

\bibitem{nuseibeh2000roadmap} B. Nuseibeh and S. Easterbrook, ``Requirements
engineering: A roadmap,'' in \emph{Proc. Conf. on the Future of Software
Engineering}, 2000, pp. 35--46.

\bibitem{bittner2014shared} E. A. C. Bittner and J. M. Leimeister, ``Creating
shared understanding in heterogeneous work groups: Why it matters and how to
achieve it,'' \emph{Journal of Management Information Systems}, vol. 31, no. 1,
pp. 111--144, 2014.

\bibitem{fuller2021blurring} R. C. Fuller and P. Kruchten, ``Blurring
boundaries: Toward the collective empathic understanding of product
requirements,'' \emph{Information and Software Technology}, vol. 140, art. no.
106670, 2021.

\bibitem{asadabadi2020ambiguity} M. R. Asadabadi, M. Saberi, O. Zwikael, and E.
Chang, ``Ambiguous requirements: A semiautomated approach to identify and clarify
ambiguity in large-scale projects,'' \emph{Computers and Industrial Engineering},
vol. 149, art. no. 106828, 2020.

\bibitem{fucci2022traceability} D. Fucci, E. Alégroth, and T. Axelsson, ``When
traceability goes awry: An industrial experience report,'' \emph{Journal of
Systems and Software}, vol. 192, art. no. 111389, 2022.

\bibitem{lim2021data} S. Lim, A. Henriksson, and J. Zdravkovic, ``Data-driven
requirements elicitation: A systematic literature review,'' \emph{SN Computer
Science}, vol. 2, art. no. 16, 2021.

\bibitem{behutiye2022quality} W. Behutiye, P. Rodríguez, M. Oivo, S. Aaramaa, J.
Partanen, and A. Abhervé, ``Towards optimal quality requirement documentation in
agile software development: A multiple case study,'' \emph{Journal of Systems and
Software}, vol. 183, art. no. 111112, 2022.

\bibitem{khan2022stakeholders} F. M. Khan, J. A. Khan, M. Assam, et al., ``A
comparative systematic analysis of stakeholder's identification methods in
requirements elicitation,'' \emph{IEEE Access}, vol. 10, pp. 30982--31011, 2022.

\bibitem{damian2006empirical} D. E. Damian and J. Chisan, ``An empirical study of
the complex relationships between requirements engineering processes and other
processes that lead to payoffs in productivity, quality, and risk management,''
\emph{IEEE Transactions on Software Engineering}, vol. 32, no. 7, pp. 433--453,
2006.

\bibitem{petersen2008systematic} K. Petersen, R. Feldt, S. Mujtaba, and M.
Mattsson, ``Systematic mapping studies in software engineering,'' in \emph{Proc.
12th Int. Conf. on Evaluation and Assessment in Software Engineering (EASE)},
2008, pp. 68--77.

\bibitem{brinkkemper1996method} S. Brinkkemper, ``Method engineering:
Engineering of information systems development methods and tools,''
\emph{Information and Software Technology}, vol. 38, no. 4, pp. 275--280, 1996.

\bibitem{henderson2010sme} B. Henderson-Sellers and J. Ralyté, ``Situational
method engineering: State-of-the-art review,'' \emph{Journal of Universal
Computer Science}, vol. 16, no. 3, pp. 424--478, 2010.

\bibitem{ralyte2001assembly} J. Ralyté and C. Rolland, ``An assembly process
model for method engineering,'' in \emph{Advanced Information Systems Engineering
(CAiSE)}, 2001, pp. 267--283.

\bibitem{mirbel2006sme} I. Mirbel and J. Ralyté, ``Situational method
engineering: Combining assembly-based and roadmap-driven approaches,''
\emph{Requirements Engineering}, vol. 11, no. 1, pp. 58--78, 2006.

\bibitem{kalus2013tailoring} G. Kalus and M. Kuhrmann, ``Criteria for software
process tailoring: A systematic review,'' in \emph{Proc. Int. Conf. on Software
and System Process (ICSSP)}, 2013, pp. 171--180.

\bibitem{campanelli2015agile} A. S. Campanelli and F. S. Parreiras, ``Agile
methods tailoring — A systematic literature review,'' \emph{Journal of Systems
and Software}, vol. 110, pp. 85--100, 2015.

\bibitem{basili1994gqm} V. R. Basili, G. Caldiera, and H. D. Rombach, ``The Goal
Question Metric approach,'' in \emph{Encyclopedia of Software Engineering}, Wiley,
1994.

\bibitem{hart1988nasa} S. G. Hart and L. E. Staveland, ``Development of NASA-TLX
(Task Load Index): Results of empirical and theoretical research,'' \emph{Advances
in Psychology}, vol. 52, pp. 139--183, 1988.

\bibitem{brooke1996sus} J. Brooke, ``SUS: A quick and dirty usability scale,'' in
\emph{Usability Evaluation in Industry}, Taylor \& Francis, 1996, pp. 189--194.

\bibitem{agredo2022thesis} V. Agredo-Delgado, \emph{A process to improve
computer-supported collaborative work through the construction, monitoring and
assistance of shared understanding in problem-solving activities}, Ph.D.
dissertation, Universidad del Cauca, Popayán, Colombia, 2022.

\bibitem{agredo2022exploratory} V. Agredo-Delgado, P. H. Ruiz, C. A. Collazos,
and F. Moreira, ``An exploratory study on the validation of THUNDERS: A process
to achieve shared understanding in problem-solving activities,''
\emph{Informatics}, vol. 9, no. 2, art. no. 39, 2022.

\bibitem{wang2025personas} Y. Wang et al., ``Who uses personas in requirements
engineering: The practitioners' perspective,'' \emph{Information and Software
Technology}, vol. 178, art. no. 107609, 2025.

\bibitem{ralph2018paradigms} P. Ralph, ``The two paradigms of software
development research,'' \emph{Science of Computer Programming}, vol. 156, pp.
68--89, 2018.

\bibitem{mattmann2015quality} I. Mattmann et al., ``The inscrutable jungle of
quality criteria — How to formulate requirements for a successful product
development,'' \emph{Procedia CIRP}, vol. 36, pp. 153--158, 2015.

\bibitem{alsanoosy2019power} T. Alsanoosy et al., ``The influence of power
distance on requirements engineering activities,'' \emph{Procedia Computer
Science}, 2019.

\bibitem{grande2024devops} R. Grande et al., ``Is it worth adopting DevOps
practices in global software engineering? Possible challenges and benefits,''
\emph{Computer Standards \& Interfaces}, vol. 87, art. no. 103767, 2024.

\bibitem{dikert2016largescale} K. Dikert, M. Paasivaara, and C. Lassenius,
``Challenges and success factors for large-scale agile transformations: A
systematic literature review,'' \emph{Journal of Systems and Software}, vol. 119,
pp. 87--108, 2016.

\bibitem{spijkman2021alignment} T. Spijkman et al., ``Alignment and granularity
of requirements and architecture in agile development: A functional
perspective,'' \emph{Information and Software Technology}, vol. 133, art. no.
106535, 2021.

\bibitem{li2024adaptability} Q. Li and F. Zeng, ``Enhancing software architecture
adaptability: A comprehensive evaluation method,'' \emph{Symmetry}, vol. 16, no.
7, art. no. 894, 2024.

\bibitem{getson2022communication} J. M. Getson et al., ``Understanding
scientists' communication challenges at the intersection of climate and
agriculture,'' \emph{PLOS ONE}, vol. 17, no. 8, art. no. e0269927, 2022.

\bibitem{lopez2021crowdre} F. López, T. Da Silva, J. A. Pino, and A. Martínez,
``CrowdRE4DF: A crowd-based requirements engineering method for developing
digital farming solutions,'' \emph{IFIP AICT}, vol. 632, pp. 310--326, 2021.

\bibitem{sulaiman2021ict} R. Sulaiman, A. Hall, T. S. V. Reddy, and K. Dorai,
``ICTs and innovation systems in agricultural development: From knowledge sharing
to knowledge co-creation,'' \emph{Agricultural Systems}, vol. 188, art. no.
103026, 2021.

\end{thebibliography}

\end{document}